\documentclass[10pt,a4paper]{amsart}
\usepackage[margin=19mm]{geometry}
\usepackage{amsmath,amssymb,mathtools}
\usepackage[scr=rsfs,cal=euler]{mathalfa}
\usepackage{xfrac}
\usepackage[unicode,hidelinks,bookmarks=false]{hyperref}
\newcommand{\ie}{i.e.\ }
\newcommand{\cf}{cf.\ }
\newcommand{\resp}{respectively\ }
\newcommand{\Subsection}{\subsection}
\providecommand{\nobreakdash}{\nobreak\hbox{-}}
\setlength{\emergencystretch}{2em}
\multlinegap0pt
\usepackage{amscd,graphics}
\usepackage{mathtools}
\usepackage{enumitem}
\usepackage{mathrsfs}
\usepackage{bm}
\usepackage{booktabs}
\usepackage{needspace}
\usepackage{tikz-cd}
  \setlength{\abovedisplayskip}{4.8pt plus 2pt minus 1pt}%
  \setlength{\belowdisplayskip}{4.8pt plus 2pt minus 1pt}%
  \setlength{\abovedisplayshortskip}{0pt plus 2pt}%
  \setlength{\belowdisplayshortskip}{4.8pt plus 2pt minus 1pt}%
\newtheorem{theorem}{Theorem}[section]
\newtheorem{proposition}[theorem]{Proposition}
\newtheorem{lemma}[theorem]{Lemma}
\newtheorem{corollary}[theorem]{Corollary}
\newtheorem{definition}[theorem]{Definition}
\newtheorem{example}[theorem]{Example}
\newtheorem{remark}[theorem]{Remark}
\newcommand{\Rep}{\operatorname{rep}}
\newcommand{\Hom}{\operatorname{Hom}}
\newcommand{\Ext}{\operatorname{Ext}}
\newcommand{\grade}{\operatorname{grade}}
\newcommand{\EssIm}{\operatorname{EssIm}}
\newcommand{\fdgrmod}{\operatorname{fdgrmod}}
\newcommand{\modcat}{\operatorname{mod}}
\newcommand{\LFgr}{\mathsf{LF}^{\mathrm{gr}}(S)}
\newcommand{\Dgr}{D_{\mathrm{gr}}}
\begin{document}
\title[Categories of Multigraded Local Cohomology Modules: Serre Fi]{Categories of Multigraded Local Cohomology Modules: Serre Filtrations and Nakayama Duality}
\author{Dancheng Lu}
\date{}
\maketitle
\noindent\emph{Source and licence.} Categories of Multigraded Local Cohomology Modules: Serre Filtrations and Nakayama Duality. Version arXiv:2609.11219v2 (CC BY 4.0). This material is distributed under the Creative Commons Attribution 4.0 International licence, http://creativecommons.org/licenses/by/4.0/, as recorded in the arXiv OAI licence metadata for this exact version and shown on the abstract page https://arxiv.org/abs/2609.11219v2. Full rights notice: licenses/arxiv-2609.11219-CC-BY-4.0.txt.\par\smallskip
\noindent\textit{Editorial scope.} Counted unit: the original \texttt{theorem} environment \texttt{thm:higher-nakayama-rank-complex} at source lines 2491--2529 together with its complete proof at lines 2531--2556; the delivered document keeps source lines 357--2557 so that every referenced label and every citation resolves inside it. Native editable source is the paper's own concatenated TeX; the AMS wrapper is editorial formatting only. No publisher class, upstream script or unrelated artwork is redistributed.
\section{Determined modules and finite-box models}\label{sec:determined}
In this section, we set up the categorical framework for the local-cohomology image categories by recalling positively and negatively $\mathbf t$-determined modules and realizing them through finite-box representation categories.

Let $\LFgr$ be the abelian category of locally finite $\mathbb Z^n$-graded
$S$-modules, i.e., modules each of whose graded components is finite-dimensional over
$\Bbbk$. For any $M\in\LFgr$, let
\[
\Dgr M=\bigoplus_{\mathbf a\in \mathbb Z^n}\Hom_\Bbbk(M_{-\mathbf a},\Bbbk)
\]
denote the graded $\Bbbk$-dual of $M$, which inherits a natural
$\mathbb Z^n$-graded $S$-module structure. The functor $\Dgr$ is the
multigraded Matlis duality; it restricts to an exact contravariant equivalence
between finitely generated and Artinian $\mathbb Z^n$-graded
$S$-modules.  This is the standard multigraded form of graded Matlis duality;
cf. \cite[Theorem~3.6.17]{BH}. We write
$D=\Hom_\Bbbk(-,\Bbbk)$ for ordinary finite-dimensional duality.

Let $\mathcal A^+$ be the full subcategory of $\LFgr$ consisting of finitely
generated $\mathbb Z^n$-graded $S$-modules $M$ such that
$M_{\mathbf a}=0$ for $\mathbf a\notin\mathbb N^n$.  For $\mathbf t\in\mathbb N^n$, let $\mathcal A^+_{\mathbf t}$ be the full
subcategory of $\mathcal A^+$ consisting of the positively
$\mathbf t$-determined modules, that is, those modules $M$ for which
multiplication by $x_j$ induces an isomorphism
\[
x_j:M_{\mathbf a}\xrightarrow{\sim}M_{\mathbf a+\varepsilon_j}
\]
for every $\mathbf a\in\mathbb N^n$ and every $j\in[n]$ such that
$a_j\ge t_j$.
The class $\mathcal A^+_{\mathbf t}$ is closed under kernels, cokernels, and
extensions, hence is abelian. Let $P_{\mathbf t}=[\mathbf0,\mathbf t]
=\{\mathbf a\in\mathbb N^n\mid \mathbf0\le\mathbf a\le\mathbf t\}$,
viewed as a subposet of $\mathbb N^n$.

\begin{proposition}[Positive finite-box equivalence]\label{prop:positive-box}
For $M\in\mathcal A^+_{\mathbf t}$, let $\pi_{\mathbf t}(M)$ be the
representation of $P_{\mathbf t}$ with
$\pi_{\mathbf t}(M)_{\mathbf a}=M_{\mathbf a}$ and structure map
$M_{\mathbf a}\to M_{\mathbf b}$ given by multiplication by
$x^{\mathbf b-\mathbf a}$ whenever $\mathbf a\le\mathbf b$.
Then restriction induces an exact equivalence
\[
\pi_{\mathbf t}:\mathcal A^+_{\mathbf t}\xrightarrow{\sim}\Rep(P_{\mathbf t}).
\]
Moreover $\mathcal A^+=\bigcup_{\mathbf t\in\mathbb N^n}\mathcal A^+_{\mathbf t}$.
\end{proposition}
\begin{proof}
This is the standard finite-box description of positively determined modules
\cite{BF,Miller}.  Explicitly, if $\mathbf a\wedge\mathbf t$ denotes the
coordinatewise minimum, the quasi-inverse sends a representation $V$ to the
module whose component in degree $\mathbf a\in\mathbb N^n$ is
$V_{\mathbf a\wedge\mathbf t}$, with identity structure maps after a
coordinate reaches $t_j$.  Every finitely generated multigraded module is positively
$\mathbf t$-determined for $\mathbf t$ sufficiently large \cite{Miller}.
\end{proof}

Following Brun--Fl\o ystad \cite[Remark~2.9]{BF}, an Artinian graded module
$N$ is \emph{negatively $\mathbf t$-determined} if
\[
N_{\mathbf a}=0\quad\text{when some }a_j\ge t_j+1,
\qquad
x_j:N_{\mathbf a}\xrightarrow{\sim}N_{\mathbf a+\varepsilon_j}
\quad(a_j\le-1).
\]
Let $\mathcal A^-_{\mathbf t}$ be the resulting abelian full exact subcategory
of Artinian graded modules and put
$\mathcal A^-=\bigcup_{\mathbf t\in\mathbb N^n}\mathcal A^-_{\mathbf t}$.

\begin{proposition}[Negative finite-box equivalence]\label{prop:negative-box}
Restriction to nonnegative degrees induces an exact equivalence
\[
\rho_{\mathbf t}:\mathcal A^-_{\mathbf t}
\xrightarrow{\sim}\Rep(P_{\mathbf t}).
\]
\end{proposition}

\begin{proof}
For $\mathbf b\in\mathbb Z^n$, put
$\mathbf b^+=(\max\{b_1,0\},\ldots,\max\{b_n,0\})$.  If
$\mathbf b\le\mathbf t$, negative determination gives a canonical isomorphism
$N_{\mathbf b}\cong N_{\mathbf b^+}$.  Hence restriction is full
and faithful.  Conversely, for $V\in\Rep(P_{\mathbf t})$ define
\[
(\sigma_{\mathbf t}V)_{\mathbf b}=
\begin{cases}
V_{\mathbf b^+},&\mathbf b\le\mathbf t,\\
0,&\text{otherwise},
\end{cases}
\]
using identity maps in negative directions and the structure maps of $V$ in
the box.  Then $\sigma_{\mathbf t}V$ is negatively $\mathbf t$-determined.
Its graded dual is generated in the finite box $[-\mathbf t,\mathbf0]$, so it
is finitely generated; graded Matlis duality therefore implies that
$\sigma_{\mathbf t}V$ is Artinian.  This gives the quasi-inverse.
\end{proof}

\begin{proposition}[Global finite-support model]\label{prop:global-fs}
Restriction to nonnegative degrees induces an exact equivalence
\[
\rho:\mathcal A^-\xrightarrow{\sim}\Rep_{\mathrm{fs}}(\mathbb N^n).
\]
\end{proposition}

\begin{proof}
Every object of $\mathcal A^-$ lies in some finite box, and every finite-support
representation is supported in one.  The preceding reconstructions are
compatible with enlarging the box.
\end{proof}

For a $\mathbb{Z}^n$-graded module $M$, we use the shift convention
$M(-\mathbf c)_{\mathbf a}=M_{\mathbf a-\mathbf c}$. By \cite[Proposition~2.8 and Remark~2.9]{BF}, if
$M\in\mathcal A^+_{\mathbf t}$ then
$H^i_{\mathfrak m}(M(-\mathbf1))\in\mathcal A^-_{\mathbf t}$.  We therefore
write
\[
\Phi_i^{\mathbf t}:\mathcal A^+_{\mathbf t}\longrightarrow
\mathcal A^-_{\mathbf t},\qquad
M\longmapsto H^i_{\mathfrak m}(M(-\mathbf1)),
\]
and define
\[
\mathcal H_i(\mathbf t)=\EssIm\Phi_i^{\mathbf t},
\qquad
\mathcal H_i=\bigcup_{\mathbf t\in\mathbb N^n}\mathcal H_i(\mathbf t).
\]
If $\mathbf t\le\mathbf u$, then
$\mathcal H_i(\mathbf t)\subseteq\mathcal H_i(\mathbf u)$.

\section{Incidence algebras, Auslander regularity, and grade}\label{sec:homological}
In this section, we collect the basic homological background needed for the first main result of the next section.

Throughout the paper, all complexes are cohomologically graded.  Thus, for a cochain complex $X$, we have $(X[r])^m=X^{m+r}$.  For a cochain complex $X$ of finite-dimensional $\Bbbk$-vector spaces, its dual $DX$ is graded by $(DX)^m=D(X^{-m})$.  A projective resolution $\cdots\to P_1\to P_0\to M\to0$ is regarded as a cochain complex with $P_r$ in degree $-r$.  Thus $\Hom(P_\bullet,A)$ is concentrated in nonnegative degrees, and applying $D$ reverses those degrees.

\subsection{The derived Nakayama description}

For a finite-dimensional $\Bbbk$-algebra $A$, we use the classical Nakayama
functor on finitely generated projective left $A$-modules in the form
$\nu_A=D\Hom_A(-,A)$, equivalently $DA\otimes_A-$; see \cite{ARS}.  Its
derived form is $DA\otimes_A^{\mathbf L}-$.  For a finite poset $P$, its
\emph{incidence algebra} $A_P$ over $\Bbbk$
is the $\Bbbk$-vector space with basis
\[
\{e_{x,y}\mid x,y\in P,\ y\le x\},
\]
whose multiplication is determined by
\[
e_{x,y}e_{z,w}=\delta_{y,z}e_{x,w}.
\]
Its identity element is $\sum_{x\in P}e_{x,x}$.  This is the convention
used by Brun--Fl\o ystad \cite{BF}.  A left $A_P$-module $M$ determines a
covariant representation of $P$ by setting $M_x=e_{x,x}M$ and letting the
map $M_y\to M_x$ for $y\le x$ be multiplication by $e_{x,y}$.
Consequently,
\[
A_P\text{-mod}\simeq\Rep(P),
\qquad
\modcat\text{-}A_P\simeq\Rep(P^{\mathrm{op}}).
\]
All module categories over finite-dimensional $\Bbbk$-algebras below
consist of finite-dimensional modules. We put
$\Lambda_{\mathbf t}=A_{P_{\mathbf t}}$ and use these equivalences for
$P=P_{\mathbf t}$ throughout.  For
$\mathbf a\in P_{\mathbf t}$, let $S_{\mathbf a}$ denote the vertex simple in
$\Rep(P_{\mathbf t})$, and let $R_{\mathbf a}$ denote the simple right
$\Lambda_{\mathbf t}$-module corresponding to the vertex simple at
$\mathbf a$ in $\Rep(P_{\mathbf t}^{\mathrm{op}})$.

For later comparison with the source and target categories, recall the
finite-box equivalences from Propositions~\ref{prop:positive-box} and
\ref{prop:negative-box}:
\[
\pi_{\mathbf t}:\mathcal A^+_{\mathbf t}\xrightarrow{\sim}\Rep(P_{\mathbf t}),
\qquad
\rho_{\mathbf t}:\mathcal A^-_{\mathbf t}\xrightarrow{\sim}\Rep(P_{\mathbf t}).
\]
Consider the order isomorphism
\[
\tau_{\mathbf t}:P_{\mathbf t}\longrightarrow P_{\mathbf t}^{\mathrm{op}},
\qquad
\tau_{\mathbf t}(\mathbf a)=\mathbf t-\mathbf a.
\]
Precomposition with $\tau_{\mathbf t}$ and its inverse gives mutually inverse
reindexing equivalences
\[
\tau_{\mathbf t}^*: \Rep(P_{\mathbf t}^{\mathrm{op}})
\longrightarrow \Rep(P_{\mathbf t}),
\qquad
\tau_{\mathbf t}^*: \Rep(P_{\mathbf t})
\longrightarrow \Rep(P_{\mathbf t}^{\mathrm{op}}),
\]
both given on vertices by
$$(\tau_{\mathbf t}^*V)_{\mathbf a}=V_{\mathbf t-\mathbf a}.$$
We use the same notation $\tau_{\mathbf t}^*$ for these two inverse functors.
Recall that
$D=\Hom_\Bbbk(-,\Bbbk)$ is the ordinary finite-dimensional $\Bbbk$-dual.  Vertexwise duality gives exact contravariant equivalences; equivalently,
viewing them as covariant functors from opposite categories,
\[
D:(\modcat\text{-}\Lambda_{\mathbf t})^{\mathrm{op}}
\xrightarrow{\sim}
\Lambda_{\mathbf t}\text{-mod},
\qquad
D:\Rep(P_{\mathbf t}^{\mathrm{op}})^{\mathrm{op}}
\xrightarrow{\sim}
\Rep(P_{\mathbf t}).
\]
This is distinct from the graded dual $\Dgr$ introduced in Section~\ref{sec:determined}.  Since $\tau_{\mathbf t}^2=\mathrm{id}$ and finite-dimensional duality commutes with reindexing, there is a natural isomorphism
\[
\tau_{\mathbf t}^*\circ D\circ\tau_{\mathbf t}^*\cong D.
\]
Thus the two order reversals occurring in the Brun--Fl\o ystad formula cancel.  In particular,
\[
D(X)_{\mathbf a}=D(X_{\mathbf a}),
\]
so $D$ preserves vertex support.

Let
\[
\mathcal T_i(\mathbf t):=\rho_{\mathbf t}(\mathcal H_i(\mathbf t))
\subseteq\Rep(P_{\mathbf t}).
\]
Since $\rho_{\mathbf t}:\mathcal A^-_{\mathbf t}\xrightarrow{\sim}\Rep(P_{\mathbf t})$
is an exact equivalence, it restricts to an exact equivalence
\[
\rho_{\mathbf t}|_{\mathcal H_i(\mathbf t)}:
\mathcal H_i(\mathbf t)\xrightarrow{\sim}\mathcal T_i(\mathbf t).
\]
Thus $\mathcal T_i(\mathbf t)$ is the representation-theoretic realization of
$\mathcal H_i(\mathbf t)$ inside the finite box.  The following statement is
the form of the Brun--Fl\o ystad Nakayama description that we shall use.

\begin{proposition}[Nakayama--Ext translation]\label{prop:naka-ext}
For $0\le i\le n$,
\[
\mathcal T_i(\mathbf t)
=
\EssIm\left(
D\Ext^{n-i}_{\Lambda_{\mathbf t}}(-,\Lambda_{\mathbf t})
:\Lambda_{\mathbf t}\text{-mod}\longrightarrow\Rep(P_{\mathbf t})
\right).
\]
\end{proposition}

\begin{proof}
Let $M\in\mathcal A^+_{\mathbf t}$ and put $X=\pi_{\mathbf t}(M).$
Choose a projective resolution $F_\bullet\to M$ in
$\mathcal A^+_{\mathbf t}$ and put $P_\bullet=\pi_{\mathbf t}(F_\bullet).$
Then $P_\bullet\to X$ is a projective resolution over
$\Lambda_{\mathbf t}$.

Brun--Fl\o ystad \cite[Proposition~3.9]{BF}, after viewing the
positively $\mathbf t$-determined complexes through the finite-box
equivalence $\pi_{\mathbf t}$, give a natural quasi-isomorphism
\[
(\tau_{\mathbf t}^*\circ D\circ\tau_{\mathbf t}^*)\!\left(
\Hom_{\Lambda_{\mathbf t}}(P_\bullet,\Lambda_{\mathbf t})
\right)
\simeq
\bigl((\pi_{\mathbf t}\circ\mathcal N_{\mathbf t}^{\mathbf1})(M)\bigr)[n].
\]
Since $\tau_{\mathbf t}^*\circ D\circ\tau_{\mathbf t}^*\cong D$, this becomes
\[
D\Hom_{\Lambda_{\mathbf t}}
(P_\bullet,\Lambda_{\mathbf t})
\simeq
\bigl((\pi_{\mathbf t}\circ\mathcal N_{\mathbf t}^{\mathbf1})(M)\bigr)[n].
\]
Here $\mathcal N_{\mathbf t}^{\mathbf1}$ denotes the first full
Nakayama functor of Brun--Fl\o ystad.  By
\cite[Section~2.3]{BF},
\[
\rho_{\mathbf t}\!\left(
H^i_{\mathfrak m}(M(-\mathbf1))
\right)
\cong
\pi_{\mathbf t}\!\left(
H^i\mathcal N_{\mathbf t}^{\mathbf1}(M)
\right).
\]
Since $\pi_{\mathbf t}$ is exact,
\[
\pi_{\mathbf t}\!\left(
H^i\mathcal N_{\mathbf t}^{\mathbf1}(M)
\right)
\cong
H^i\!\left(
(\pi_{\mathbf t}\circ\mathcal N_{\mathbf t}^{\mathbf1})(M)
\right).
\]
Therefore
\[
\begin{aligned}
\rho_{\mathbf t}\!\left(
H^i_{\mathfrak m}(M(-\mathbf1))
\right)
&\cong
H^{i-n}\!\left(
\bigl((\pi_{\mathbf t}\circ\mathcal N_{\mathbf t}^{\mathbf1})(M)\bigr)[n]
\right)\\
&\cong
H^{i-n}\!\left(
D\Hom_{\Lambda_{\mathbf t}}
(P_\bullet,\Lambda_{\mathbf t})
\right)\\
&\cong
D\Ext_{\Lambda_{\mathbf t}}^{\,n-i}
(X,\Lambda_{\mathbf t}).
\end{aligned}
\]
\end{proof}

\subsection{Auslander regular algebras}

We recall only the properties needed below.  Let $A$ be a finite-dimensional $\Bbbk$-algebra, and consider a minimal injective coresolution of the right regular module
\[
0\longrightarrow A_A\longrightarrow I^0\longrightarrow I^1\longrightarrow I^2\longrightarrow\cdots.
\]

\begin{definition}
For $k\ge1$, the algebra $A$ is \emph{$k$-Gorenstein} if
\[
\operatorname{pd}_A I^r\le r
\qquad (0\le r<k).
\]
It satisfies the \emph{Auslander condition} if it is $k$-Gorenstein for every $k$.  It is \emph{Auslander regular} if it satisfies the Auslander condition and has finite global dimension.
\end{definition}

Our convention for $k$-Gorenstein algebras and the Auslander condition agrees with that used by Iyama--Marczinzik \cite{IM}.  For a right $A$-module $X$, define
$$
\grade_{A^{\mathrm{op}}} X
=
\inf\{r\ge0\mid \Ext^r_{A^{\mathrm{op}}}(X,A)\ne0\},
$$
with the convention $\inf\varnothing=\infty$; in particular,
$\grade_{A^{\mathrm{op}}}0=\infty$.
By Auslander's theorem \cite[Theorem~3.7]{FGR}, the Auslander condition is equivalent to the following condition: for every left $A$-module $M$, every $q\ge1$, and every right submodule
$$
X\subseteq\Ext_A^q(M,A),
$$
one has $\grade_{A^{\mathrm{op}}} X\ge q.$


\subsection{A realization lemma for Ext}

The following elementary consequence of the Auslander condition is the homological mechanism behind the Serre filtration.

\begin{lemma}\label{lem:realization}
Let $A$ be a finite-dimensional algebra satisfying the Auslander condition and let $q\ge1$.  Then
\[
\EssIm\left(
\Ext_A^q(-,A):A\text{-}\modcat\longrightarrow\modcat\text{-}A
\right)
=
\{X\in\modcat\text{-}A\mid \grade_{A^{\mathrm{op}}} X\ge q\}.
\]
\end{lemma}

\begin{proof}
If $X=\Ext_A^q(M,A)$, the Auslander condition gives $\grade_{A^{\mathrm{op}}} X\ge q$.

Conversely, let $X$ be a right $A$-module with $\grade_{A^{\mathrm{op}}} X\ge q$ and choose a projective resolution
\[
\cdots\longrightarrow P_q\longrightarrow P_{q-1}\longrightarrow\cdots\longrightarrow P_0\longrightarrow X\longrightarrow0.
\]
Write $(-)^*=\Hom_{A^{\mathrm{op}}}(-,A)$.  Since
\[
\Ext^r_{A^{\mathrm{op}}}(X,A)=0
\qquad (0\le r<q),
\]
dualizing the first $q$ steps yields an exact sequence
\[
0\longrightarrow P_0^*\longrightarrow P_1^*\longrightarrow\cdots\longrightarrow P_q^*\longrightarrow M\longrightarrow0,
\]
where
\[
M=\operatorname{coker}(P_{q-1}^*\longrightarrow P_q^*).
\]
This is a projective resolution of the left $A$-module $M$ of length $q$.
Here $P_q^*$ occurs in degree $0$ and $P_0^*$ in degree $q$.  Applying
$\Hom_A(-,A)$ and using $P_r^{**}\cong P_r$, the last cohomology group is
therefore the cokernel of $P_1\to P_0$.  Hence
\[
\Ext_A^q(M,A)
\cong
\operatorname{coker}(P_1\longrightarrow P_0)
\cong X.
\]
\end{proof}

For $q\ge1$ put
\[
\mathcal G_q(A)
=
\{X\in\modcat\text{-}A\mid \grade_{A^{\mathrm{op}}} X\ge q\}.
\]

\begin{proposition}\label{prop:grade-serre}
If $A$ satisfies the Auslander condition, then $\mathcal G_q(A)$ is a Serre subcategory of $\modcat\text{-}A$ for every $q\ge1$.
\end{proposition}

\begin{proof}
Extension closure follows from the long exact Ext sequence.  If $U\subseteq X$ and $X\in\mathcal G_q(A)$, Lemma~\ref{lem:realization} writes $X\cong\Ext_A^q(M,A)$, and the Auslander condition gives $\grade_{A^{\mathrm{op}}} U\ge q$.  Thus $\mathcal G_q(A)$ is closed under submodules.  Finally, if
\[
0\longrightarrow U\longrightarrow X\longrightarrow Q\longrightarrow0
\]
has $U,X\in\mathcal G_q(A)$, the long exact Ext sequence gives
\[
\Ext^r_{A^{\mathrm{op}}}(Q,A)=0
\qquad (0\le r<q),
\]
so $Q\in\mathcal G_q(A)$.
\end{proof}

\section{Local cohomology categories are abelian except in top degree}\label{sec:finitebox}

We now specialize the preceding homological results to the incidence algebra
$\Lambda_{\mathbf t}$ of the finite box $P_{\mathbf t}$.  For $i<n$, the
local-cohomology image categories are Serre, hence abelian.  In top degree the
behavior changes: the image is the category of second cosyzygies, which is
extension-closed but, for every nontrivial box, non-abelian.  We first establish
these statements in a finite box and then pass to the global form.

The poset $P_{\mathbf t}$ is a finite distributive lattice.  Iyama and
Marczinzik prove that its incidence algebra is Auslander regular
\cite[Theorem~3.2]{IM}, and compute the grade of every simple module
\cite[Corollary~3.3]{IM}.  Their incidence-algebra convention is opposite to
the Brun--Fl\o ystad convention used here.  Accordingly, right modules over
$\Lambda_{\mathbf t}$ correspond to the left modules in the convention of
\cite{IM}, or equivalently to representations of the opposite poset.  Thus
their formula translates, for the right simple $R_{\mathbf a}$ over
$\Lambda_{\mathbf t}$, to
\[
\grade_{\Lambda_{\mathbf t}^{\mathrm{op}}}R_{\mathbf a}
=
\#\{\text{elements of $P_{\mathbf t}$ covered by $\mathbf a$}\}
=
\ell(\mathbf a).
\]
Indeed, $\mathbf a$ covers one element in coordinate $j$ exactly when
$a_j>0$.  This is the convention needed in the proof below.

\subsection{The layers below the top}

We now prove the finite-box classification announced in the Introduction.

\begin{theorem}[Intrinsic finite-box description]\label{thm:finite-intrinsic}
Let $0\le i<n$ and put $q=n-i$.  Then
\[
\boxed{
\mathcal T_i(\mathbf t)
=
\left\{
V\in\Rep(P_{\mathbf t})\mid
V_{\mathbf a}=0\text{ whenever }\ell(\mathbf a)<q
\right\}
=
\Rep(U_q(\mathbf t)).
}
\]
Consequently $\mathcal T_i(\mathbf t)$ is a Serre subcategory of $\Rep(P_{\mathbf t})$, and
$\rho_{\mathbf t}$ restricts to an exact equivalence
\[
\rho_{\mathbf t}|_{\mathcal H_i(\mathbf t)}:
\mathcal H_i(\mathbf t)\xrightarrow{\sim}\Rep(U_q(\mathbf t)).
\]
\end{theorem}

\begin{proof}
Since $q\ge1$ and $\Lambda_{\mathbf t}$ is Auslander regular, the algebra $\Lambda_{\mathbf t}$ satisfies the Auslander condition.  Therefore Proposition~\ref{prop:naka-ext}, Lemma~\ref{lem:realization}, and Proposition~\ref{prop:grade-serre} give
\[
\mathcal T_i(\mathbf t)
=
D\bigl(\mathcal G_q(\Lambda_{\mathbf t})\bigr).
\]
Since $D$ is an exact duality and $\mathcal G_q(\Lambda_{\mathbf t})$ is a Serre subcategory of $\modcat\text{-}\Lambda_{\mathbf t}$, it follows already that $\mathcal T_i(\mathbf t)$ is a Serre subcategory of $\Rep(P_{\mathbf t})$.

It remains to identify this Serre subcategory explicitly.  Since $\modcat\text{-}\Lambda_{\mathbf t}$ has finite length, $\mathcal G_q(\Lambda_{\mathbf t})$ is determined by the simple modules it contains.  The right simple $R_{\mathbf a}$ belongs to $\mathcal G_q(\Lambda_{\mathbf t})$ precisely when
\[
\grade_{\Lambda_{\mathbf t}^{\mathrm{op}}}R_{\mathbf a}=\ell(\mathbf a)\ge q.
\]
The duality $D$ preserves vertex support.  Hence the corresponding left representations have no composition factors at vertices satisfying $\ell(\mathbf a)<q$.  Since the dimension vector is additive in short exact sequences and the vertex simple at $\mathbf a$ has dimension vector concentrated at $\mathbf a$, this condition is equivalent to
\[
V_{\mathbf a}=0\qquad\text{for }\ell(\mathbf a)<q.
\]
Since $\ell(\mathbf a)\le\ell(\mathbf b)$ whenever
$\mathbf a\le\mathbf b$, the subposet $U_q(\mathbf t)$ is an upper order
ideal of $P_{\mathbf t}$.  Hence extension by zero identifies
$\Rep(U_q(\mathbf t))$ with the full subcategory of
$\Rep(P_{\mathbf t})$ consisting of representations $V$ such that
$\operatorname{supp}V\subseteq U_q(\mathbf t)$.  This gives the asserted
identification.
\end{proof}

\begin{corollary}\label{cor:finite-filtration}
We have
\[
\mathcal T_0(\mathbf t)
\subseteq\mathcal T_1(\mathbf t)
\subseteq\cdots\subseteq\mathcal T_{n-1}(\mathbf t).
\]
For $1\le i<n$, the new simple objects appearing from
$\mathcal T_{i-1}(\mathbf t)$ to $\mathcal T_i(\mathbf t)$ are the vertex
simples $S_{\mathbf a}$ with
\[
\ell(\mathbf a)=n-i.
\]
If $\mathbf t\ge\mathbf1$, all inclusions are strict.
\end{corollary}

The following proposition gives a syzygy explanation for the filtration in Corollary~\ref{cor:finite-filtration}. We first record the following elementary observation. Let $N$ be a finitely generated $\mathbb N^n$-graded $S$-module. If $N$ is positively $\mathbf t$-determined, then all its minimal homogeneous generators have multidegrees $\le \mathbf t$. Conversely, if $N$ is torsion-free and generated in multidegrees $\le \mathbf t$, then $N$ is positively $\mathbf t$-determined.

\begin{proposition}\label{prop:syzygy-filtration}
Let $M\in\mathcal A^+_{\mathbf t}$, let $P\twoheadrightarrow M$ be a minimal graded free cover, and put
$$
\Omega M=\ker(P\to M).
$$
Then $\Omega M\in\mathcal A^+_{\mathbf t}$. Moreover, for
$0\le i\le n-2$,
$$
H^i_{\mathfrak m}(M(-\mathbf1))
\cong
H^{i+1}_{\mathfrak m}(\Omega M(-\mathbf1)).
$$
\end{proposition}

\begin{proof}
Since $M$ is positively $\mathbf t$-determined, all its minimal homogeneous generators have multidegrees $\le \mathbf t$. Hence its minimal graded free cover has the form
$$
P=
\bigoplus_{\mathbf b\le \mathbf t}
S(-\mathbf b)^{\beta_{0,\mathbf b}(M)}.
$$

Thus $P$ is torsion-free and generated in multidegrees $\le \mathbf t$. By the elementary observation above, $P$ is positively $\mathbf t$-determined. Hence $P,M\in\mathcal A^+_{\mathbf t}.$


Since $\mathcal A^+_{\mathbf t}$ is an abelian full subcategory of $\mathcal A^+$, it is closed under kernels. Therefore
$$
\Omega M=\ker(P\to M)\in\mathcal A^+_{\mathbf t}.
$$

Now shift the short exact sequence
$$
0\longrightarrow\Omega M\longrightarrow P\longrightarrow M\longrightarrow0
$$
by $-\mathbf1$. Since $P$ is free, $H^r_{\mathfrak m}(P(-\mathbf1))=0$ for $0\le r<n$.
The associated long exact sequence in local cohomology therefore yields the desired isomorphism.
\end{proof}




\begin{remark}\label{rem:dimension-shift-stops}
Proposition~\ref{prop:syzygy-filtration} not only explains why $\mathcal H_i(\mathbf t)\subseteq\mathcal H_{i+1}(\mathbf t)$ for $0\le i\le n-2$, but also explains intrinsically why the Serre chain stops at
$\mathcal H_{n-1}(\mathbf t)$.  At $i=n-1$ the long exact sequence contains
\[
0\longrightarrow H^{n-1}_{\mathfrak m}(M(-\mathbf1))
\longrightarrow H^n_{\mathfrak m}(\Omega M(-\mathbf1))
\longrightarrow H^n_{\mathfrak m}(P(-\mathbf1)),
\]
and the last term is generally nonzero.  Thus dimension shifting no longer
gives an isomorphism with the top local cohomology category.
\end{remark}

\subsection{The top-degree category and second cosyzygies}

For $i=n$, Proposition~\ref{prop:naka-ext} gives
\[
\mathcal T_n(\mathbf t)
=
\EssIm\left(
D\Hom_{\Lambda_{\mathbf t}}(-,\Lambda_{\mathbf t})
:\Lambda_{\mathbf t}\text{-mod}\longrightarrow\Rep(P_{\mathbf t})
\right).
\]
The grade-realization lemma no longer applies because the relevant Ext degree
is zero.  On the right-module side the replacement is the category of second
syzygies; after applying $D$, the resulting objects in
$\Rep(P_{\mathbf t})$ are second cosyzygies.

For a finite-dimensional $\Bbbk$-algebra $A$, let
$\operatorname{Syz}_2(\modcat\text{-}A)$ denote the full subcategory of
$\modcat\text{-}A$ consisting of all modules $X$ for which there exists an
exact sequence
$$
0\longrightarrow X\longrightarrow P_1\longrightarrow P_0
\longrightarrow C\longrightarrow0
$$
with $P_0$ and $P_1$ projective.  Thus nonminimal projective resolutions are
allowed; in particular, projective modules belong to
$\operatorname{Syz}_2(\modcat\text{-}A)$.

\begin{lemma}\label{lem:dual-syzygy}
For every finite-dimensional $\Bbbk$-algebra $A$,
$$
\EssIm\!\left(
\Hom_A(-,A):A\text{-}\modcat\longrightarrow\modcat\text{-}A
\right)
=
\operatorname{Syz}_2(\modcat\text{-}A).
$$
\end{lemma}

\begin{proof}
Let $M\in A\text{-}\modcat$, put $M^*=\Hom_A(M,A)$, and choose a projective
presentation
$$
P_1\longrightarrow P_0\longrightarrow M\longrightarrow0.
$$
Applying $(-)^*=\Hom_A(-,A)$ gives an exact sequence
$$
0\longrightarrow M^*\longrightarrow P_0^*\longrightarrow P_1^*.
$$
Since $P_0^*$ and $P_1^*$ are projective right $A$-modules, this shows that
$M^*\in\operatorname{Syz}_2(\modcat\text{-}A)$.

Conversely, let $X\in\operatorname{Syz}_2(\modcat\text{-}A)$.  By definition,
there is an exact sequence
$$
0\longrightarrow X\longrightarrow Q_1\xrightarrow{d}Q_0
\longrightarrow C\longrightarrow0
$$
with $Q_0$ and $Q_1$ projective right $A$-modules.  Put
$$
M=\operatorname{coker}(d^*:Q_0^*\longrightarrow Q_1^*).
$$
Then $Q_0^*\to Q_1^*\to M\to0$ is a projective presentation of the left
$A$-module $M$.  Dualizing it and using the natural isomorphisms
$Q_i^{**}\cong Q_i$ gives
$$
M^*\cong\ker(d)=X.
$$
Thus $X$ belongs to the essential image of $\Hom_A(-,A)$.
\end{proof}

We shall use the following result of Auslander--Reiten \cite{ARsyzygy}.  If an Artin algebra $A$ is $k$-Gorenstein for every $k$, that is, if a minimal injective resolution
\[
0\longrightarrow A_A\longrightarrow I^0\longrightarrow I^1\longrightarrow\cdots
\]
satisfies $\operatorname{pd}_A I^j\le j$ for every $j\ge0$, then
the class of $r$-th syzygies (allowing nonminimal projective resolutions) is extension-closed for every $r\ge1$.
Since $\Lambda_{\mathbf t}$ is Auslander regular, it is $k$-Gorenstein for every $k$, so in particular
$\operatorname{Syz}_2(\modcat\text{-}\Lambda_{\mathbf t})$ is extension-closed.

For $\mathbf a\in P_{\mathbf t}$ let $I^{\mathbf t}_{\mathbf a}$ denote the principal lower-interval representation
\[
I^{\mathbf t}_{\mathbf a}(\mathbf b)=
\begin{cases}
\Bbbk,&\mathbf b\le\mathbf a,\\
0,&\text{otherwise},
\end{cases}
\]
with identity maps between nonzero components.  These are precisely the indecomposable injective objects of $\Rep(P_{\mathbf t})$.

Dually, we call an object $V$ of an abelian category a \emph{second cosyzygy}
if there exists an exact sequence
$$
0\longrightarrow W\longrightarrow I^0\longrightarrow I^1
\longrightarrow V\longrightarrow0
$$
with $I^0$ and $I^1$ injective.

\begin{theorem}[Explicit description of the top finite-box category]\label{thm:top-finite}
For $V\in\Rep(P_{\mathbf t})$, the following are equivalent:
\begin{enumerate}[label=\rm(\arabic*)]
\item $V\in\mathcal T_n(\mathbf t)$;
\item there is an exact sequence
\[
0\longrightarrow W\longrightarrow I^0\longrightarrow I^1\longrightarrow V\longrightarrow0
\]
with $I^0$ and $I^1$ injective in $\Rep(P_{\mathbf t})$;
\item $V$ is the cokernel of a morphism between finite direct sums of principal lower-interval representations,
\[
\bigoplus_{\mathbf a\in P_{\mathbf t}}
   (I^{\mathbf t}_{\mathbf a})^{m_{\mathbf a}}
\longrightarrow
\bigoplus_{\mathbf a\in P_{\mathbf t}}
   (I^{\mathbf t}_{\mathbf a})^{n_{\mathbf a}}.
\]
\end{enumerate}
Thus $\mathcal T_n(\mathbf t)$ is the category of second cosyzygies in $\Rep(P_{\mathbf t})$.  In particular, it contains every injective representation and is extension-closed in $\Rep(P_{\mathbf t})$.
\end{theorem}

\begin{proof}
By Lemma~\ref{lem:dual-syzygy},
\[
\mathcal T_n(\mathbf t)
=
D\bigl(\operatorname{Syz}_2(\modcat\text{-}\Lambda_{\mathbf t})\bigr).
\]
An object $X$ belongs to $\operatorname{Syz}_2(\modcat\text{-}\Lambda_{\mathbf t})$ precisely when it occurs in an exact sequence
\[
0\longrightarrow X\longrightarrow Q_1\longrightarrow Q_0
\longrightarrow C\longrightarrow0
\]
with $Q_0,Q_1$ projective.  Since $D$ is an exact contravariant duality sending projectives to injectives, applying $D$ gives an exact sequence
\[
0\longrightarrow D(C)\longrightarrow D(Q_0)\longrightarrow D(Q_1)
\longrightarrow D(X)\longrightarrow0.
\]
Thus $D(X)$ is a second cosyzygy.  Conversely, applying $D$ to a second
cosyzygy sequence for $V$ gives such a second-syzygy sequence, with
$X=D(V)$.  This proves the equivalence of (1) and (2).  Since every injective
in the finite-length category $\Rep(P_{\mathbf t})$ is a finite direct sum of
the $I^{\mathbf t}_{\mathbf a}$, (2) and (3) are equivalent.

Finally, $\operatorname{Syz}_2(\modcat\text{-}\Lambda_{\mathbf t})$ is extension-closed by the Auslander--Reiten result stated above, and exact duality preserves extension closure.  Hence $\mathcal T_n(\mathbf t)$ is extension-closed.
\end{proof}

The top layer is not part of the Serre filtration.

\begin{proposition}\label{prop:top-disjoint}
For every $0\le i<n$,
\[
\mathcal T_i(\mathbf t)\cap\mathcal T_n(\mathbf t)=0.
\]
\end{proposition}

\begin{proof}
By Theorem~\ref{thm:finite-intrinsic} it is enough to show
\[
\operatorname{Syz}_2(\modcat\text{-}\Lambda_{\mathbf t})
\cap
\mathcal G_1(\Lambda_{\mathbf t})=0.
\]
If $0\ne X\in\operatorname{Syz}_2(\modcat\text{-}\Lambda_{\mathbf t})$, then $X$ embeds in a finitely generated projective right module $P$.  Since $P$ is a direct summand of $\Lambda_{\mathbf t}^m$ for some $m$, there is a nonzero map $X\to\Lambda_{\mathbf t}^m$, hence a nonzero coordinate map $X\to\Lambda_{\mathbf t}$.  Thus
\[
\Hom_{\Lambda_{\mathbf t}^{\mathrm{op}}}(X,\Lambda_{\mathbf t})\ne0,
\]
so $\grade_{\Lambda_{\mathbf t}^{\mathrm{op}}}X=0$.  Hence $X\notin\mathcal G_1(\Lambda_{\mathbf t})$.
\end{proof}

\begin{proposition}[Non-abelianity in top degree]\label{prop:top-nonabelian}
Assume that $P_{\mathbf t}$ has a nonzero vertex, equivalently
$\mathbf t\ne\mathbf0$.  Then $\mathcal T_n(\mathbf t)$, and hence
$\mathcal H_n(\mathbf t)$, is not an abelian category.
\end{proposition}

\begin{proof}
Choose $\mathbf a\ne\mathbf0$ and let
\[
\pi:I^{\mathbf t}_{\mathbf a}\longrightarrow
I^{\mathbf t}_{\mathbf0}=S_{\mathbf0}
\]
be the canonical epimorphism which is the identity at $\mathbf0$ and zero
elsewhere.  Both terms lie in $\mathcal T_n(\mathbf t)$.  Its kernel $K$ in
$\Rep(P_{\mathbf t})$ is nonzero and satisfies $K_{\mathbf0}=0$; hence
$K\in\mathcal T_{n-1}(\mathbf t)$ by
Theorem~\ref{thm:finite-intrinsic}, and therefore
$K\notin\mathcal T_n(\mathbf t)$ by Proposition~\ref{prop:top-disjoint}.
Thus $\mathcal T_n(\mathbf t)$ is not closed under kernels computed in
$\Rep(P_{\mathbf t})$.  To show that it is not even an abelian category in
its own right, we argue intrinsically as follows.

For every $V\in\mathcal T_n(\mathbf t)$, Theorem~\ref{thm:top-finite}(3)
shows that each structure map $V_{\mathbf0}\to V_{\mathbf b}$ is surjective:
this is true for every principal injective and is preserved by finite direct
sums and cokernels.  Since $K_{\mathbf0}=0$, naturality gives
$\Hom(V,K)=0$.  Thus $\pi$ is a monomorphism in
$\mathcal T_n(\mathbf t)$; it is also an epimorphism there, since it is one in
$\Rep(P_{\mathbf t})$, but it is not an isomorphism.  An abelian category is
balanced, so $\mathcal T_n(\mathbf t)$ is not abelian.  The assertion for
$\mathcal H_n(\mathbf t)$ follows from $\rho_{\mathbf t}$.
\end{proof}

Combining Proposition~\ref{prop:top-disjoint} with Theorem~\ref{thm:finite-intrinsic} for $i=n-1$ gives
\[
0\ne V\in\mathcal T_n(\mathbf t)
\quad\Longrightarrow\quad
V_{\mathbf0}\ne0.
\]
Moreover, the only vertex simple in the top category is $S_{\mathbf0}$.  Indeed, $S_{\mathbf0}=I^{\mathbf t}_{\mathbf0}$ is injective and hence belongs to $\mathcal T_n(\mathbf t)$ by Theorem~\ref{thm:top-finite}, whereas every $S_{\mathbf a}$ with $\mathbf a\ne\mathbf0$ belongs to $\mathcal T_{n-1}(\mathbf t)$ and is therefore excluded by Proposition~\ref{prop:top-disjoint}.

For an object $X$ of an additive category, write $\operatorname{add}X$ for
the full subcategory of direct summands of finite direct sums of copies of
$X$.  For full additive subcategories $\mathcal X,\mathcal Y$ of an abelian
category, write $\mathcal X\bullet\mathcal Y$ for the additive closure of
objects $E$ occurring in a short exact sequence
\[
0\longrightarrow X\longrightarrow E\longrightarrow Y\longrightarrow0,
\qquad X\in\mathcal X,\quad Y\in\mathcal Y.
\]
We use the standard convention that a pair $(\mathcal T,\mathcal F)$ of full
subcategories is a \emph{torsion pair} if
$\Hom(\mathcal T,\mathcal F)=0$ and every object $E$ occurs in a short exact
sequence $0\to T\to E\to F\to0$ with $T\in\mathcal T$ and
$F\in\mathcal F$.  It is \emph{hereditary} if $\mathcal T$ is closed under
subobjects, equivalently if $\mathcal T$ is a Serre subcategory.

The order-ideal recollement developed in
Section~\ref{sec:general-fs-theory} will recover all successive-layer torsion
pairs uniformly.  The passage from $\mathcal H_{n-1}(\mathbf t)$ to the exceptional top category
$\mathcal H_n(\mathbf t)$ behaves differently.

\begin{proposition}[Top--penultimate extension decomposition]\label{prop:top-penultimate}
Inside $\Rep(P_{\mathbf t})$ one has
\[
\Rep(P_{\mathbf t})
=
\mathcal T_{n-1}(\mathbf t)\bullet\operatorname{add}S_{\mathbf0}
=
\mathcal T_{n-1}(\mathbf t)\bullet\mathcal T_n(\mathbf t).
\]
Moreover
\[
\bigl(\mathcal T_{n-1}(\mathbf t),\operatorname{add}S_{\mathbf0}\bigr)
\]
is a hereditary torsion pair in $\Rep(P_{\mathbf t})$.  In general
$\bigl(\mathcal T_{n-1}(\mathbf t),\mathcal T_n(\mathbf t)\bigr)$ is
not a torsion pair.
\end{proposition}

\begin{proof}
By Theorem~\ref{thm:finite-intrinsic},
\[
\mathcal T_{n-1}(\mathbf t)
=
\{V\in\Rep(P_{\mathbf t})\mid V_{\mathbf0}=0\}.
\]
For $V\in\Rep(P_{\mathbf t})$, define a subrepresentation $V^{>0}$ by
\[
(V^{>0})_{\mathbf0}=0,
\qquad
(V^{>0})_{\mathbf a}=V_{\mathbf a}
\quad(\mathbf a\ne\mathbf0).
\]
This is a subrepresentation because $\mathbf0$ is the minimal element of
$P_{\mathbf t}$.  Hence
$V^{>0}\in\mathcal T_{n-1}(\mathbf t)$ and
\[
V/V^{>0}\cong
S_{\mathbf0}^{\,\dim_\Bbbk V_{\mathbf0}}.
\]
Since $S_{\mathbf0}\in\mathcal T_n(\mathbf t)$, the two asserted
extension equalities follow.

If $X\in\mathcal T_{n-1}(\mathbf t)$, then $X_{\mathbf0}=0$, and hence
$\Hom(X,S_{\mathbf0})=0$.  The displayed short exact sequence therefore
shows that
$\bigl(\mathcal T_{n-1}(\mathbf t),\operatorname{add}S_{\mathbf0}\bigr)$
is a torsion pair.  It is hereditary because
$\mathcal T_{n-1}(\mathbf t)$ is a Serre subcategory.

Finally, if $\mathbf a\ne\mathbf0$, then
$S_{\mathbf a}\in\mathcal T_{n-1}(\mathbf t)$ and the principal
lower-interval representation
$I^{\mathbf t}_{\mathbf a}$ belongs to $\mathcal T_n(\mathbf t)$,
while
\[
\Hom(S_{\mathbf a},I^{\mathbf t}_{\mathbf a})\ne0.
\]
Thus the required Hom-orthogonality for
$\bigl(\mathcal T_{n-1}(\mathbf t),\mathcal T_n(\mathbf t)\bigr)$
fails whenever the box has a nonzero vertex.
\end{proof}


\subsection{Global representation model and the top category}\label{sec:global}

The finite-box theorem globalizes particularly cleanly through Proposition~\ref{prop:global-fs}.  For $q\ge0$ set
\[
U_q=\{\mathbf a\in\mathbb N^n\mid \ell(\mathbf a)\ge q\}.
\]

\begin{theorem}[Global representation model]\label{thm:global-model}
Let $0\le i<n$ and put $q=n-i$.  Then the equivalence
\[
\rho:\mathcal A^-\xrightarrow{\sim}\Rep_{\mathrm{fs}}(\mathbb N^n)
\]
restricts to an exact equivalence
\[
\mathcal H_i\xrightarrow{\sim}\Rep_{\mathrm{fs}}(U_q).
\]
Consequently $\mathcal H_i$ is a Serre subcategory of $\mathcal A^-$ and
\[
\mathcal H_0\subsetneq\mathcal H_1\subsetneq\cdots\subsetneq\mathcal H_{n-1}.
\]
\end{theorem}

\begin{proof}
Take $N\in\mathcal A^-$ and choose $\mathbf t$ such that $N\in\mathcal A^-_{\mathbf t}$.  By Theorem~\ref{thm:finite-intrinsic}, $N$ belongs to $\mathcal H_i(\mathbf t)$ if and only if its nonnegative restriction vanishes at every vertex $\mathbf a$ with $\ell(\mathbf a)<q$.  This is exactly the condition that $\rho(N)$ be supported on $U_q$.  Since support conditions define Serre subcategories in a functor category, the assertion follows.  Strictness is witnessed by the vertex simples with $\ell(\mathbf a)=q$.
\end{proof}

\begin{corollary}\label{cor:global-top}
The category $\mathcal H_n$ is extension-closed in $\mathcal A^-$, and
\[
\mathcal H_i\cap\mathcal H_n=0
\qquad(0\le i<n).
\]
If $n\ge1$, then $\mathcal H_n$ is not an abelian category.
\end{corollary}

\begin{proof}
Given a short exact sequence in $\mathcal A^-$ with end terms in
$\mathcal H_n$, choose $\mathbf t$ large enough that all three terms lie in
$\mathcal A^-_{\mathbf t}$ and apply Theorem~\ref{thm:top-finite}.  The
intersection statement follows similarly from
Proposition~\ref{prop:top-disjoint} in a common finite box.

For the last assertion, use the canonical epimorphism
$I_{\varepsilon_1}\to S_{\mathbf0}$.  Given $V\in\mathcal H_n$, choose a
finite box containing both $V$ and $\varepsilon_1$; the proof of
Proposition~\ref{prop:top-nonabelian} in that box shows that this map is a
monomorphism when tested against $V$.  Hence it is both monic and epic in
$\mathcal H_n$ but not an isomorphism, so $\mathcal H_n$ is not abelian.
\end{proof}

\begin{corollary}\label{cor:global-top-extension}
At the exceptional top step,
\[
\mathcal A^-=
\mathcal H_{n-1}\bullet\operatorname{add}S_{\mathbf0}
=
\mathcal H_{n-1}\bullet\mathcal H_n,
\]
whereas $(\mathcal H_{n-1},\mathcal H_n)$ is not a torsion pair.
\end{corollary}

\begin{proof}
The assertion follows from Proposition~\ref{prop:top-penultimate} inside a
common finite box.  The same nonzero morphism
$S_{\mathbf a}\to I_{\mathbf a}$, for $\mathbf a\ne\mathbf0$, witnesses the
failure of Hom-orthogonality globally.
\end{proof}

\section{Finite-support representations of inward-finite posets}
\label{sec:general-fs-theory}

The aim of this section is to develop a general framework for representations
of posets needed in the subsequent sections.  The posets may be infinite, but
all representations have finite-dimensional vertex spaces and finite support.
Since the infinite posets occurring later are inward finite, we formulate the
one-sided theory in that direction.  The outward-finite theory is obtained by
passing to opposite posets and will not be stated separately.  For finite
posets, of course, both orientations are available.

\subsection{Kan extensions and recollements}
\label{subsec:recollement-background}
We use the standard formalism of abelian recollements as in
\cite{FP,Psaroudakis,PV} and of triangulated recollements as in \cite{BBD}.
Thus an abelian recollement
consists of adjoint triples
\[
i^*\dashv i_*\dashv i^!,
\qquad
j_!\dashv j^*\dashv j_*,
\]
with $i_*,j_!,j_*$ fully faithful and $\operatorname{Im}i_*=\ker j^*$.
We shall repeatedly use the following standard consequences.

\begin{proposition}[TTF and quotient consequences]\label{prop:recollement-ttf}
Let
\[
\begin{tikzcd}[column sep=5.5em]
\mathcal A'
  \arrow[r, "i_*" description]
&
\mathcal A
  \arrow[l, bend right=20, "i^*"']
  \arrow[l, bend left=20, "i^!"]
  \arrow[r, "j^*" description]
&
\mathcal A''
  \arrow[l, bend left=32, "j_!"']
  \arrow[l, bend right=32, "j_*"]
\end{tikzcd}
\]
be an abelian recollement.
\begin{enumerate}[label=\rm(\arabic*)]
\item The triple
\[
\bigl(\ker i^*,\,\operatorname{Im}i_*,\,\ker i^!\bigr)
\]
is a TTF triple in $\mathcal A$; that is,
$(\ker i^*,\operatorname{Im}i_*)$ and
$(\operatorname{Im}i_*,\ker i^!)$ are torsion pairs.
\item The exact functor $j^*$ induces an equivalence
\[
\mathcal A/\operatorname{Im}i_*
\xrightarrow{\sim}
\mathcal A''.
\]
\item If $i^*$ is exact, then $\ker i^*$ is a Serre subcategory,
$(\ker i^*,\operatorname{Im}i_*)$ is a hereditary torsion pair, and
$i^*$ induces an equivalence
\[
\mathcal A/\ker i^*
\xrightarrow{\sim}
\mathcal A'.
\]
\item Dually, if $i^!$ is exact, then $\ker i^!$ is a Serre subcategory and
$i^!$ induces an equivalence
\[
\mathcal A/\ker i^!
\xrightarrow{\sim}
\mathcal A'.
\]
\end{enumerate}
\end{proposition}

These are standard consequences of an abelian recollement.  For the
correspondence between abelian recollements, strong TTF triples, and
bilocalising Serre subcategories, see \cite[Theorem~4.3]{PV}; for the standard
recollement localization properties, see \cite[Section~2]{Psaroudakis}.
Statements~(3) and~(4) also follow from Gabriel's localization theorem for
an exact functor admitting a fully faithful adjoint; see \cite{Gabriel}.

For an arbitrary poset $P$, call $P$ \emph{inward finite} if the lower set
$\{y\in P\mid y\le x\}$ is finite for every $x\in P$.  Dually, $P$ is
\emph{outward finite} if the upper set $\{y\in P\mid x\le y\}$ is finite for
every $x\in P$.  The latter notion will be used only to describe the
order-dual situation and to indicate when an opposite orientation may fail
globally.

An injective map $f:C\to P$ of posets is called a \emph{full inclusion} if
$x\le y$ in $C$ if and only if $f(x)\le f(y)$ in $P$.  Equivalently, it is a
fully faithful functor between the corresponding categories.  When
$C\subseteq P$ has the order induced from $P$, we call $C$ a subposet of $P$;
its inclusion is full.

Let $f:C\hookrightarrow P$ be a full inclusion.  We write $f^*$ for
restriction along $f$, and denote its left and right Kan extensions by
$f_!$ and $f_*$, respectively.  Thus
\[
f_!\dashv f^*\dashv f_*.
\]
For a $\Bbbk$-representation $V$ of $C$ and $p\in P$, the two Kan extensions are
given pointwise by
\[
(f_!V)_p=\operatorname*{colim}_{\substack{c\in C\\c\le p}}V_c,
\qquad
(f_*V)_p=\operatorname*{lim}_{\substack{c\in C\\p\le c}}V_c,
\]
respectively.  We use the convention that the colimit or limit of an empty
diagram is $0$.  If $C$ is an upper order ideal, then $f_!$ is extension by
zero; if $C$ is a lower order ideal, then $f_*$ is extension by zero.

\begin{lemma}\label{lem:fs-kan-support}
Let $f:C\hookrightarrow P$ be a full inclusion and
$V\in\Rep_{\mathrm{fs}}(C)$.  Then
\[
\operatorname{supp}(f_*V)
\subseteq\bigcup_{c\in\operatorname{supp}V}\{p\in P\mid p\le c\}.
\]
Consequently, if $P$ is inward finite, then $f_*$ preserves finite support.
\end{lemma}

\begin{proof}
If no point of $\operatorname{supp}V$ lies above $p$, then every nonzero term
is absent from the diagram defining $(f_*V)_p$, so this limit is zero.  The
containment follows, and inward finiteness makes the right-hand side a finite
union of finite principal lower sets.
\end{proof}

\begin{proposition}[Inward-finite order-ideal recollement]
\label{prop:fs-order-recollement}
Let $Q$ be an upper order ideal of an inward-finite poset $P$ and put
$R=P\setminus Q$, so that $R$ is a lower order ideal.  Write
$u:Q\hookrightarrow P$ and $v:R\hookrightarrow P$ for the inclusions.  Then
$u_*$ preserves finite support and there is an abelian recollement
\[
\begin{tikzcd}[column sep=5.5em]
\Rep_{\mathrm{fs}}(R)
  \arrow[r, "v_*" description]
& \Rep_{\mathrm{fs}}(P)
  \arrow[l, bend right=20, "v^*"']
  \arrow[l, bend left=20, "i^!"]
  \arrow[r, "u^*" description]
& \Rep_{\mathrm{fs}}(Q)
  \arrow[l, bend left=20, "u_!"']
  \arrow[l, bend right=20, "u_*"]
\end{tikzcd}
\]
Here $u_!$ and $v_*$ are extension by zero, and in the standard recollement
notation
\[
i^*=v^*,\qquad i_*=v_*,\qquad
j_!=u_!,\qquad j^*=u^*,\qquad j_*=u_*.
\]
The remaining functor is
\[
i^!M
=
v^*\operatorname{Ker}\!\left(M\longrightarrow u_*u^*M\right).
\]
Equivalently, for $r\in R$,
\[
(i^!M)_r
=
\bigcap_{\substack{q\in Q\\ r\le q}}
\operatorname{Ker}(M_r\longrightarrow M_q),
\]
where the maps are the structure maps of $M$; if there is no $q\in Q$ with
$r\le q$, the intersection is understood to be $M_r$.  The structure maps
of $i^!M$ are induced by those of $M$.


\end{proposition}

\begin{proof}
By Lemma~\ref{lem:fs-kan-support}, $u_*$ preserves finite support, so
\[
u_!\dashv u^*\dashv u_*
\]
restricts to the finite-support representation categories.  Moreover,
$\ker u^*=\operatorname{Im}v_*$.  Since $v_*$ is extension by zero, its
left adjoint is the restriction $v^*$.  A morphism $v_*W\to M$ is determined
on $R$, and compatibility with every relation $r\le q$, where $r\in R$ and
$q\in Q$, forces its image at $r$ to lie in
\[
\bigcap_{\substack{q\in Q\\ r\le q}}
\operatorname{Ker}(M_r\longrightarrow M_q).
\]
Thus $v_*\dashv i^!$ and the displayed formula for $i^!$ follows.  Hence the
six functors give the recollement.  If $P$ is finite, the same argument in the
order-dual poset gives the opposite orientation.
\end{proof}

\begin{corollary}[Order-ideal torsion and quotients]
\label{cor:order-ideal-consequences}
Under the hypotheses of Proposition~\ref{prop:fs-order-recollement},
\[
\bigl(\Rep_{\mathrm{fs}}(Q),\Rep_{\mathrm{fs}}(R)\bigr)
\]
is a hereditary torsion pair in $\Rep_{\mathrm{fs}}(P)$.  Moreover,
restriction induces equivalences
\[
\Rep_{\mathrm{fs}}(P)/\Rep_{\mathrm{fs}}(R)
\xrightarrow{\sim}\Rep_{\mathrm{fs}}(Q),
\qquad
\Rep_{\mathrm{fs}}(P)/\Rep_{\mathrm{fs}}(Q)
\xrightarrow{\sim}\Rep_{\mathrm{fs}}(R).
\]
\end{corollary}

\begin{proof}
In the recollement of Proposition~\ref{prop:fs-order-recollement},
$i^*=v^*$ is the exact restriction functor,
\[
\ker i^*=\Rep_{\mathrm{fs}}(Q),
\qquad
\operatorname{Im}i_*=\Rep_{\mathrm{fs}}(R).
\]
Thus Proposition~\ref{prop:recollement-ttf}(1),(3) gives the hereditary
torsion pair and the second quotient equivalence, while
Proposition~\ref{prop:recollement-ttf}(2) gives the first.
\end{proof}


A descending chain of upper order ideals may therefore be treated
successively.  Iterating Corollary~\ref{cor:order-ideal-consequences} gives a
canonical functorial torsion filtration whose successive factors are supported
on the corresponding strata.  We shall use this observation below for the
support-rank filtration of the local-cohomology categories.

\begin{remark}\label{rem:finite-complementary-recollements}
If $P$ is finite, every Kan extension has finite support, so the order-dual
construction gives the complementary recollement with the two side categories
interchanged.  The displayed recollement is also the standard idempotent
recollement of $A_P$ associated with $e_Q=\sum_{q\in Q}e_q$, since
$e_QA_Pe_Q\simeq A_Q$ and $A_P/A_Pe_QA_P\simeq A_R$.
\end{remark}

The inward-finite hypothesis also provides the boundedness needed below.
Indeed, for $V\in\Rep_{\mathrm{fs}}(P)$ its lower closure
\[
P_V=\bigcup_{p\in\operatorname{supp}V}\{x\in P\mid x\le p\}
\]
is a finite lower ideal.  Resolve $V|_{P_V}$ injectively in $\Rep(P_V)$ and
extend the resolution by zero.  Since extension by zero along a lower-ideal
inclusion is right adjoint to exact restriction, it preserves injectives.
Thus every finite-support representation has finite injective dimension; the
same holds for every subposet of $P$.

\begin{corollary}[Bounded-derived order-ideal recollement]
\label{cor:fs-derived-recollement}
Under the hypotheses of Proposition~\ref{prop:fs-order-recollement}, the
abelian recollement lifts to the triangulated recollement
\[
\begin{tikzcd}[column sep=4.2em]
D^b\!\left(\Rep_{\mathrm{fs}}(R)\right)
  \arrow[r, "v_*" description]
& D^b\!\left(\Rep_{\mathrm{fs}}(P)\right)
  \arrow[l, bend right=20, "v^*"']
  \arrow[l, bend left=20, "\mathbf R i^!"]
  \arrow[r, "u^*" description]
& D^b\!\left(\Rep_{\mathrm{fs}}(Q)\right)
  \arrow[l, bend left=20, "u_!"']
  \arrow[l, bend right=20, "\mathbf R u_*"]
\end{tikzcd}
\]
Here the exact functors $v^*,v_*,u_!$, and $u^*$ are applied degreewise,
whereas the left-exact right adjoints $i^!$ and $u_*$ are replaced by their
total right derived functors $\mathbf R i^!$ and $\mathbf R u_*$,
respectively.  If $P$ is finite, then both recollement orientations lift to
bounded derived categories.
\end{corollary}

\begin{proof}
The preceding argument with the finite lower ideal gives enough injectives, with
bounded injective resolutions, in all three finite-support categories.  Hence
$\mathbf R i^!$ and $\mathbf R u_*$ are defined on the bounded derived
categories and preserve boundedness.  Since $v^*,v_*,u_!$, and $u^*$ are
exact, they act degreewise on bounded derived categories and retain their
adjunctions.  Standard derived adjunction also gives
$v_*\dashv\mathbf R i^!$ and $u^*\dashv\mathbf R u_*$.  The identities
$v^*v_*\cong\mathrm{id}$ and $u^*u_!\cong\mathrm{id}$ imply that $v_*$ and
$u_!$ remain fully faithful.

It remains to identify the kernel of $u^*$.  For
$X^\bullet\in D^b(\Rep_{\mathrm{fs}}(P))$, one has $u^*X^\bullet\simeq0$
if and only if every $H^m(X^\bullet)$ is supported on $R$.  In this case the
unit morphism
\[
X^\bullet\longrightarrow v_*v^*X^\bullet
\]
is a quasi-isomorphism, since it is the identity on $R$ and has acyclic
components on $Q$.  Thus
$\ker u^*=\operatorname{Im}v_*$ on the bounded derived categories, and the six
functors form the displayed triangulated recollement.  This is also the
homological-embedding mechanism appearing in general derived-lifting criteria;
compare \cite[Theorem~7.2]{Psaroudakis}, where enough projectives and
injectives are assumed.  If $P$ is finite, the order-dual argument gives the
complementary orientation.
\end{proof}

\begin{remark}\label{rem:finite-kan-not-exact}
Finiteness of $P$ does not, however, imply that the right Kan extension $u_*$
is exact; it only guarantees preservation of finite support.  Passing to the
complementary recollement orientation does not change the exactness of a given
functor, but replaces the relevant nontrivial adjoints by those belonging to
the other orientation.  In the finite-box support-rank recollement below, the
two stratum Kan extensions are additionally exact by
Lemma~\ref{lem:stratum-kan-exact}; consequently only one of the six functors
in Corollary~\ref{cor:derived-horizontal-recollement} needs to be derived.
\end{remark}

\subsection{Right Serre functors and Kan exchange}

For a set $X$, write $\Bbbk[X]$ for its $\Bbbk$-linearization, that is, the
free $\Bbbk$-vector space with basis $X$.  For $p\in P$, let
\[
P_p:=\Bbbk[P(p,-)]\in\Rep(P)
\]
denote the $\Bbbk$-linearized covariant representable functor.  Its support is
the principal upper set $\{x\in P\mid p\le x\}$, so $P_p$ need not belong to
$\Rep_{\mathrm{fs}}(P)$.  The contravariant representable
$\Bbbk[P(-,p)]$ belongs to $\Rep(P^{\mathrm{op}})$, and let
\[
I_p:=D\!\bigl(\Bbbk[P(-,p)]\bigr)\in\Rep(P)
\]
be its vertexwise dual.  For every $V\in\Rep(P)$, the dual form of the
Yoneda lemma gives a natural isomorphism
\[
\Hom_{\Rep(P)}(V,I_p)\cong D(V_p).
\]
Since evaluation at $p$ and finite-dimensional $\Bbbk$-duality are exact,
$I_p$ is injective in $\Rep(P)$.  Its support is the principal lower set
$\{x\in P\mid x\le p\}$; hence inward finiteness ensures that
$I_p\in\Rep_{\mathrm{fs}}(P)$.  Since $\Rep_{\mathrm{fs}}(P)$ is a full
exact subcategory of $\Rep(P)$, $I_p$ is injective in
$\Rep_{\mathrm{fs}}(P)$ as well.  In fact, if $S_p$ denotes the vertex
simple at $p$, then the canonical monomorphism $S_p\hookrightarrow I_p$ is
essential, so $I_p$ is the injective envelope of $S_p$ in
$\Rep_{\mathrm{fs}}(P)$.

Put
\[
\mathcal P(P):=\operatorname{add}\{P_p\mid p\in P\},
\qquad
\mathcal I(P):=\operatorname{add}\{I_p\mid p\in P\}.
\]
Thus $\mathcal P(P)$ is the category of finitely generated projective objects
of $\Rep(P)$, and, when $P$ is inward finite,
$\mathcal I(P)\subseteq\Rep_{\mathrm{fs}}(P)$.  For
$Q\in\mathcal P(P)$ define
\begin{equation}\label{eq:projective-nakayama}
(\nu_P^0 Q)_x
:=
D\Hom_{\Rep(P)}(Q,P_x)
\qquad(x\in P).
\end{equation}
If $x\le y$, the natural map $P_y\to P_x$ induces, after applying
$\Hom_{\Rep(P)}(Q,-)$ and then $D$, the structure map
$(\nu_P^0Q)_x\to(\nu_P^0Q)_y$.  Hence \eqref{eq:projective-nakayama} defines
an additive functor
\[
\nu_P^0:\mathcal P(P)\longrightarrow\mathcal I(P).
\]
For $p,x\in P$, the Yoneda lemma gives a natural isomorphism
\[
\Hom_{\Rep(P)}(P_p,P_x)
\cong
(P_x)_p
=
\Bbbk[P(x,p)].
\]
Hence
\[
(\nu_P^0P_p)_x
=
D\Hom_{\Rep(P)}(P_p,P_x)
\cong
D\Bbbk[P(x,p)]
=
(I_p)_x.
\]
These isomorphisms are natural in $x$, and therefore assemble to a natural
isomorphism of representations
\[
\nu_P^0(P_p)\cong I_p.
\]

For finite $P$, finite global dimension of the incidence algebra allows
$\nu_P^0$ to extend, via bounded projective resolutions, to the derived
Nakayama functor on $D^b(\Rep(P))$; in this case it coincides with the Serre
functor.  For possibly infinite $P$, inward finiteness keeps
$I_p=\nu_P^0(P_p)$ finitely supported,
while finite $P_p$-resolutions allow passage to the bounded derived
finite-support category.  This motivates the following notion.

We call an inward-finite poset $P$ \emph{right-Serre finite} if every object
of $\Rep_{\mathrm{fs}}(P)$ admits, in $\Rep(P)$, a finite resolution by
objects of $\mathcal P(P)$, equivalently by finite direct sums of the
representables $P_p$.  The terms of such a resolution need not have finite
support.

\begin{lemma}
\label{lem:right-serre-model}
Let $P$ be inward finite.  Then $P$ is right-Serre finite if and only if every
vertex simple $S_p$ admits a finite resolution in $\Rep(P)$ by objects of
$\mathcal P(P)$.  If these equivalent conditions hold, define
\[
K^b_{\mathrm{fs}}(\mathcal P(P))
=
\left\{
Q^\bullet\in K^b(\mathcal P(P))\ \middle|\
H^r(Q^\bullet)\in\Rep_{\mathrm{fs}}(P)\text{ for all }r
\right\}.
\]
Then the canonical functor
$K^b(\mathcal P(P))\to D^b(\Rep(P))$ identifies
$K^b_{\mathrm{fs}}(\mathcal P(P))$ with the essential image of
\[
D^b(\Rep_{\mathrm{fs}}(P))\longrightarrow D^b(\Rep(P)).
\]
In particular, there is a natural triangle equivalence
\begin{equation}\label{eq:bounded-projective-model}
D^b(\Rep_{\mathrm{fs}}(P))
\simeq
K^b_{\mathrm{fs}}(\mathcal P(P)).
\end{equation}
\end{lemma}

\begin{proof}
Every finite-support representation has finite length, and its composition
factors are vertex simples.  Hence the ``only if'' direction is immediate,
while the converse follows by induction on length and the horseshoe lemma.

For the derived statement, inward finiteness implies that the principal
injectives $I_p$ have finite support.  Consequently
$\Rep_{\mathrm{fs}}(P)$ has enough injectives which are injective also in
$\Rep(P)$: explicitly, one applies the finite-lower-ideal construction given
before Corollary~\ref{cor:fs-derived-recollement}.  Every finite-support object
therefore has a finite resolution by such injectives.  It follows that the
inclusion
$\Rep_{\mathrm{fs}}(P)\hookrightarrow\Rep(P)$ induces a fully faithful
functor on bounded derived categories, since morphisms on either side may be
computed using the same finite injective resolutions.  Bounded complexes of objects of
$\mathcal P(P)$ are $K$-projective, so
$K^b(\mathcal P(P))\to D^b(\Rep(P))$ is fully faithful.  Right-Serre
finiteness allows one to resolve the finitely many terms of any bounded
complex in $\Rep_{\mathrm{fs}}(P)$ and totalize, giving a bounded complex in
$\mathcal P(P)$.  Conversely, if $Q^\bullet\in K^b(\mathcal P(P))$ has
finite-support cohomology, the standard truncation triangles express its
image in $D^b(\Rep(P))$ as an iterated extension of its finite-support
cohomology objects.  Since the fully faithful image of
$D^b(\Rep_{\mathrm{fs}}(P))$ is triangulated, the image of $Q^\bullet$ lies
in it.  This proves \eqref{eq:bounded-projective-model}.
\end{proof}

\begin{lemma}[Lower-ideal inheritance]
\label{lem:right-serre-lower-ideal}
Let $P$ be an inward-finite right-Serre finite poset and let
$j:Q\hookrightarrow P$ be a lower order ideal.  Then $Q$ is right-Serre
finite.
\end{lemma}

\begin{proof}
The lower ideal $Q$ is inward finite.  By
Lemma~\ref{lem:right-serre-model}, it is enough to resolve the vertex simples.
Fix $q\in Q$ and regard $S_q$ as a representation of $P$ by extension by
zero.  Choose a finite resolution of $S_q$ in $\Rep(P)$ by objects of
$\mathcal P(P)$ and restrict it to $Q$.  Restriction is exact, and for every
$p\in P$ one has
\[
j^*P_p^P\cong
\begin{cases}
P_p^Q,&p\in Q,\\
0,&p\notin Q.
\end{cases}
\]
Indeed, if $p\notin Q$ and $p\le q'\in Q$, the lower-ideal property would
force $p\in Q$.  Hence the restricted complex is a finite resolution of
$S_q$ by objects of $\mathcal P(Q)$, and the assertion follows from
Lemma~\ref{lem:right-serre-model}.
\end{proof}

Assume now that $P$ is right-Serre finite.  Under
\eqref{eq:bounded-projective-model}, define
\[
\mathbb S_P^R:
D^b(\Rep_{\mathrm{fs}}(P))
\longrightarrow
D^b(\Rep_{\mathrm{fs}}(P))
\]
by applying $\nu_P^0$ term by term: if $Q^\bullet\in
K^b_{\mathrm{fs}}(\mathcal P(P))$ represents $X$, then
\begin{equation}\label{eq:right-serre-definition}
\mathbb S_P^R(X):=\nu_P^0(Q^\bullet).
\end{equation}
The right-hand side is a bounded complex of objects of $\mathcal I(P)$, so it
lies in $D^b(\Rep_{\mathrm{fs}}(P))$.

\begin{proposition}[Existence of a right Serre functor]
\label{prop:right-serre-existence}
If $P$ is right-Serre finite, then $\mathbb S_P^R$ defined by
\eqref{eq:right-serre-definition} is a right Serre functor: for
$X,Y\in D^b(\Rep_{\mathrm{fs}}(P))$ there are bifunctorial isomorphisms
\[
D\Hom(X,Y)\cong\Hom(Y,\mathbb S_P^R X).
\]
If $P$ is finite, then $\mathcal P(P)=\operatorname{proj}\Rep(P)$, and
the finite global dimension of the incidence algebra of $P$ yields, via
bounded projective resolutions, an equivalence
\[
D^b(\Rep(P))\simeq K^b(\mathcal P(P)).
\]
Moreover, $\mathbb S_P^R$ agrees with the derived Nakayama functor, which is
the Serre autoequivalence of $D^b(\Rep(P))$.
\end{proposition}

\begin{proof}
For $Q\in\mathcal P(P)$ and $Y\in\Rep_{\mathrm{fs}}(P)$, Yoneda and
finite-dimensional duality give, first for $Q=P_p$ and then by additivity,
\[
D\Hom_{\Rep(P)}(Q,Y)
\cong
\Hom_{\Rep(P)}(Y,\nu_P^0Q).
\]
For a bounded projective complex $Q^\bullet$ representing $X$, these
isomorphisms are compatible with the differentials and therefore give the
corresponding duality of Hom-complexes.  Since $Q^\bullet$ is $K$-projective
and $\nu_P^0(Q^\bullet)$ is a bounded complex of injectives, taking
zeroth cohomology yields
\[
D\Hom_{D^b(\Rep_{\mathrm{fs}}(P))}(X,Y)
\cong
\Hom_{D^b(\Rep_{\mathrm{fs}}(P))}(Y,\mathbb S_P^R X).
\]
Naturality is inherited from \eqref{eq:projective-nakayama}.  If $P$ is
finite, every $P_p$ has finite support and the incidence algebra of $P$ has
finite global dimension, so the final statement is the standard
projective-to-injective description of the derived Nakayama functor, which in
this finite case is the Serre functor.
\end{proof}

\begin{remark}\label{rem:inward-not-serre}
Inward finiteness and right-Serre finiteness control different sides of the
homological picture.  The former makes the $I_p$ finite-support injectives;
the latter supplies bounded resolutions by the generally non-finite-support
projectives $P_p$.  Inward finiteness alone does not imply right-Serre
finiteness.  For example, let
\[
P=\{0,a_1,a_2,\ldots\},
\qquad
0<a_i\quad(i\ge1),
\]
with the $a_i$ pairwise incomparable.  Then $P$ is inward finite, but the
kernel of $P_0\twoheadrightarrow S_0$ is
$\bigoplus_{i\ge1}P_{a_i}$, which is not finitely generated.  Thus $S_0$
is not finitely presented: indeed, if it admitted a finite presentation by
objects of $\mathcal P(P)$, Schanuel's lemma would imply that the kernel of
the displayed finite projective cover is finitely generated, a contradiction.
In particular, $S_0$ has no finite resolution by objects of
$\mathcal P(P)$, and $P$ is not right-Serre finite.
\end{remark}

\begin{theorem}[Right-Serre exchange]\label{thm:right-serre-exchange}
Let $f:Q\hookrightarrow P$ be a full inclusion of right-Serre finite posets.
Assume that the derived Kan extensions restrict to adjoints
\[
\mathbf L f_!,\mathbf R f_*:
D^b(\Rep_{\mathrm{fs}}(Q))\longrightarrow
D^b(\Rep_{\mathrm{fs}}(P))
\]
of the exact restriction $f^*$.  Then there is a canonical natural
isomorphism
\begin{equation}\label{eq:general-right-serre-exchange}
\boxed{\mathbb S_P^R\circ\mathbf L f_!
\simeq\mathbf R f_*\circ\mathbb S_Q^R.}
\end{equation}
\end{theorem}

\begin{proof}
For $X\in D^b(\Rep_{\mathrm{fs}}(Q))$ and
$Y\in D^b(\Rep_{\mathrm{fs}}(P))$, right Serre duality and the derived
adjunctions give
\begin{align*}
\Hom(Y,(\mathbb S_P^R\circ\mathbf L f_!)(X))
&\cong D\Hom(\mathbf L f_!X,Y)\\
&\cong D\Hom(X,f^*Y)\\
&\cong\Hom(f^*Y,\mathbb S_Q^R X)\\
&\cong\Hom(Y,(\mathbf R f_*\circ\mathbb S_Q^R)(X)).
\end{align*}
Yoneda gives \eqref{eq:general-right-serre-exchange}.
\end{proof}

\section{Application I: The Finite-Box Case}
\label{sec:finite-box-recollement}

We now apply the general theory of Section~\ref{sec:general-fs-theory} to the
finite-box categories $\mathcal H_i(\mathbf t)$.  Since their representation
models $U_q(\mathbf t)$ are finite, both Kan extensions preserve finite support,
so the general theory already provides the relevant recollements and Serre
exchange.  The finite-box setting, however, allows us to go substantially
further.  We compute the horizontal Kan extensions and recollement functors
explicitly, construct a functorial rank-layer resolution comparing the two
sections $j_!$ and $j_*$, and use Nakayama--Serre duality to transform this
resolution into a costandard rank complex.  Thus this section is not merely a
specialization of Section~\ref{sec:general-fs-theory}; the finiteness of the box
makes possible the explicit homological constructions developed below.

\subsection{The two-parameter Serre stratification}

The finite-box categories vary in two independent directions: the
cohomological index and the determining box.  By
Proposition~\ref{prop:negative-box} and
Theorem~\ref{thm:finite-intrinsic}, under the global equivalence
$\rho:\mathcal A^-\xrightarrow{\sim}\Rep_{\mathrm{fs}}(\mathbb N^n)$,
$\mathcal H_i(\mathbf t)$, for $0\le i<n$, is the Serre subcategory supported
on $U_{n-i}(\mathbf t)$.  This common realization gives the following
two-parameter description.

\begin{theorem}[Two-parameter Serre stratification]\label{thm:two-parameter}
Let $0\le i,j<n$ and $\mathbf s,\mathbf t\in\mathbb N^n$, and put
$p=n-i$ and $q=n-j$.
\begin{enumerate}[label=\rm(\arabic*)]
\item If $i\le j$ and $\mathbf s\le\mathbf t$, then
$\mathcal H_i(\mathbf s)\subseteq\mathcal H_j(\mathbf t)$ is a Serre
inclusion.
\item Writing
$(\mathbf s\wedge\mathbf t)_r=\min\{s_r,t_r\}$, one has
\[
\mathcal H_i(\mathbf s)\cap\mathcal H_j(\mathbf t)
=\mathcal H_{\min\{i,j\}}(\mathbf s\wedge\mathbf t).
\]
\item If $i\le j$ and $\mathbf s\le\mathbf t$, then
\[
\mathcal H_j(\mathbf t)/\mathcal H_i(\mathbf s)
\simeq
\Rep\bigl(U_q(\mathbf t)\setminus U_p(\mathbf s)\bigr),
\]
where the complement has the order induced from $U_q(\mathbf t)$.
Explicitly, its elements are the $\mathbf a\le\mathbf t$ such that
$\ell(\mathbf a)\ge q$ and either $\mathbf a\nleq\mathbf s$ or
$\ell(\mathbf a)<p$.
\end{enumerate}
\end{theorem}

\begin{proof}
If $i\le j$ and $\mathbf s\le\mathbf t$, then $p\ge q$ and
$U_p(\mathbf s)\subseteq U_q(\mathbf t)$.  This subposet is convex: if
$\mathbf a,\mathbf c\in U_p(\mathbf s)$ and
$\mathbf a\le\mathbf b\le\mathbf c$, then
$\mathbf b\le\mathbf s$ and $\ell(\mathbf b)\ge\ell(\mathbf a)\ge p$.
Thus representations supported on $U_p(\mathbf s)$ form a Serre subcategory
of $\Rep(U_q(\mathbf t))$.  Restriction to the complementary full subposet is
exact, has this Serre subcategory as its kernel, and has a fully faithful
right Kan extension.  Gabriel localization therefore proves (1) and (3)
\cite{Gabriel}.

For (2), the common support realization gives
\[
U_p(\mathbf s)\cap U_q(\mathbf t)
=U_{\max\{p,q\}}(\mathbf s\wedge\mathbf t),
\]
which corresponds to
$\mathcal H_{\min\{i,j\}}(\mathbf s\wedge\mathbf t)$.
\end{proof}

On a finite poset the terms ``standard'' and ``costandard'' below are
descriptive: they mean objects obtained by left and right Kan extension,
respectively, and do not assert a highest-weight structure.

\subsection{The horizontal recollement}\label{subsec:horizontal-recollement}

We now study the order-ideal decomposition underlying two adjacent
cohomological layers.  Its recollement both recovers the horizontal quotient
and records how that quotient sits inside $\Rep(U_q(\mathbf t))$.  Since
$U_q(\mathbf t)$ is finite and
$U_q(\mathbf t)=U_{q+1}(\mathbf t)\sqcup L_q(\mathbf t)$ is an upper/lower
order-ideal decomposition, Proposition~\ref{prop:fs-order-recollement}
applies; the point of the results below is to compute the two horizontal
sections explicitly.

For the explicit Kan-extension calculations in the following lemma and
proposition, fix $1\le i\le n-1$ and put $q=n-i$.  For $q\le r\le n$ set
\[
L_r(\mathbf t)=\{\mathbf a\in P_{\mathbf t}\mid\ell(\mathbf a)=r\},
\qquad
j_r:L_r(\mathbf t)\hookrightarrow U_q(\mathbf t),
\]
and write $j=j_q$.  We also write
\[
\kappa:U_{q+1}(\mathbf t)\hookrightarrow U_q(\mathbf t)
\]
for the inclusion of the upper order ideal.  We use the same symbols
$j_{r,!}$ and $j_{r,*}$ for the left and right Kan extensions on
representation categories; their pointwise formulas were recalled above.
Since $U_{q+1}(\mathbf t)$ is an upper order ideal, the left Kan extension
$\kappa_!$ is extension by zero.  In the recollement notation below we write
$\iota_*:=\kappa_!$ for this embedding, and denote its left and right
adjoints by $\iota^*$ and $\iota^!$, respectively.  Thus the symbol
$\iota_*$ is recollement notation and does not denote the right Kan
extension $\kappa_*$.

For $\mathbf a\in P_{\mathbf t}$ and $F\subseteq\operatorname{supp}(\mathbf a)$, define $\mathbf a_F\in \mathbb N^n$ by
\[
(\mathbf a_F)_k=
\begin{cases}
a_k,&k\in F,\\
0,&k\notin F.
\end{cases}
\]

\begin{lemma}\label{lem:stratum-kan-exact}
For every $q\le r\le n$, both Kan extension functors
\[
j_{r,!},j_{r,*}:\Rep(L_r(\mathbf t))\longrightarrow\Rep(U_q(\mathbf t))
\]
are exact.
\end{lemma}

\begin{proof}
Fix $\mathbf a\in U_q(\mathbf t)$ and put
$A=\operatorname{supp}(\mathbf a)$.  By the pointwise formula for left Kan
extension,
\[
(j_{r,!}W)_{\mathbf a}
=
\operatorname*{colim}_{\substack{\mathbf b\in L_r(\mathbf t)\\
\mathbf b\le\mathbf a}} W_{\mathbf b}.
\]
The indexing poset is the disjoint union of the subposets consisting of those
$\mathbf b$ with fixed support
$B=\operatorname{supp}(\mathbf b)\subseteq A$, $|B|=r$.  For each such $B$,
the corresponding subposet has terminal object $\mathbf a_B$, since every
$\mathbf b\le\mathbf a$ with $\operatorname{supp}(\mathbf b)=B$ satisfies
$\mathbf b\le\mathbf a_B$.  Hence
\[
(j_{r,!}W)_{\mathbf a}
\cong
\bigoplus_{\substack{B\subseteq A\\|B|=r}} W_{\mathbf a_B},
\]
so $(j_{r,!}W)_{\mathbf a}$ is a finite direct sum of vertex evaluations of
$W$.

For the right Kan extension, the pointwise formula gives
\[
(j_{r,*}W)_{\mathbf a}
=
\operatorname*{lim}_{\substack{\mathbf b\in L_r(\mathbf t)\\
\mathbf a\le\mathbf b}} W_{\mathbf b}.
\]
If $|A|>r$, the indexing poset is empty.  If $|A|\le r$, it is the disjoint
union of the subposets with fixed support $B\supseteq A$, $|B|=r$.  The
subposet indexed by $B$ has an initial object, obtained from $\mathbf a$ by
putting the coordinates in $B\setminus A$ equal to $1$.  Thus
$(j_{r,*}W)_{\mathbf a}$ is a finite product of vertex evaluations of $W$.
Finite direct sums and finite products of finite-dimensional $\Bbbk$-vector spaces
are exact, so both Kan extensions are exact.
\end{proof}

\begin{proposition}\label{prop:horizontal-embeddings}
Restriction along $j$ is the quotient functor in the abelian recollement
\[
\begin{tikzcd}[column sep=4.8em]
\Rep(U_{q+1}(\mathbf t))
  \arrow[r, "\iota_*" description]
& \Rep(U_q(\mathbf t))
  \arrow[l, bend right=20, "\iota^*"']
  \arrow[l, bend left=20, "\iota^!"]
  \arrow[r, "j^*" description]
& \Rep(L_q(\mathbf t))
  \arrow[l, bend left=20, "j_!"']
  \arrow[l, bend right=20, "j_*"]
\end{tikzcd}
\]
Here $U_{q+1}(\mathbf t)$ is an upper order ideal of $U_q(\mathbf t)$, and
$\iota_*=\kappa_!$ is extension by zero.  Its right adjoint is restriction,
\[
(\iota^!M)_{\mathbf a}=M_{\mathbf a}
\qquad(\mathbf a\in U_{q+1}(\mathbf t)),
\]
whereas its left adjoint is given pointwise by
\[
(\iota^*M)_{\mathbf a}
=
M_{\mathbf a}\Big/
\sum_{\substack{\mathbf b\in L_q(\mathbf t)\\\mathbf b\le\mathbf a}}
\operatorname{Im}\bigl(M_{\mathbf b}\longrightarrow M_{\mathbf a}\bigr),
\qquad \mathbf a\in U_{q+1}(\mathbf t).
\]
The two adjoint sections $j_!$ and $j_*$ are exact and fully faithful.  Explicitly, if $V\in\Rep(L_q(\mathbf t))$, $\mathbf a\in U_q(\mathbf t)$, and $A=\operatorname{supp}(\mathbf a)$, then
\[
(j_!V)_{\mathbf a}
=
\bigoplus_{\substack{F\subseteq A\\|F|=q}}V_{\mathbf a_F},
\qquad
(j_*V)_{\mathbf a}
=
\begin{cases}
V_{\mathbf a},&|A|=q,\\
0,&|A|>q.
\end{cases}
\]
There is a canonical natural epimorphism
\[
\epsilon_V:j_!V\longrightarrow j_*V
\]
which is the identity on $L_q(\mathbf t)$ and zero on $U_{q+1}(\mathbf t)$.
The associated TTF triple contains the hereditary torsion pair
\[
\bigl(\operatorname{Im}\iota_*,\operatorname{Im}j_*\bigr)
=
\bigl(\Rep(U_{q+1}(\mathbf t)),j_*\Rep(L_q(\mathbf t))\bigr)
\quad\text{in }\Rep(U_q(\mathbf t)),
\]
where both side categories in the second expression are identified with their
extension-by-zero images.
\end{proposition}

\begin{proof}
For $\mathbf a\in U_q(\mathbf t)$, put $A=\operatorname{supp}(\mathbf a)$.  The indexing poset for the pointwise left Kan extension consists of the elements $\mathbf b\in L_q(\mathbf t)$ with $\mathbf b\le\mathbf a$.  It is the disjoint union of the subposets having fixed support $F\subseteq A$, $|F|=q$, and the subposet indexed by $F$ has terminal object $\mathbf a_F$.  Its colimit is therefore the displayed direct sum.

For the right Kan extension, the indexing poset consists of the elements $\mathbf b\in L_q(\mathbf t)$ satisfying $\mathbf a\le\mathbf b$.  If $|A|>q$, this category is empty.  If $|A|=q$, it has the initial object $\mathbf a$, so its limit is $V_{\mathbf a}$.  These formulas show directly that both Kan extensions are exact and that $j^*j_!\cong\mathrm{id}\cong j^*j_*$.  Hence both are fully faithful.

The existence of this complementary recollement follows from
Proposition~\ref{prop:fs-order-recollement} and
Remark~\ref{rem:finite-complementary-recollements}, since
$U_q(\mathbf t)$ is finite.  It remains only to identify the adjoints to
$\iota_*$.  A morphism
$\iota_*N\to M$ is uniquely determined by its restriction to
$U_{q+1}(\mathbf t)$, giving
$\iota^!M=M|_{U_{q+1}(\mathbf t)}$.  A morphism $M\to\iota_*N$ must vanish
on $L_q(\mathbf t)$, so naturality forces its component at
$\mathbf a\in U_{q+1}(\mathbf t)$ to annihilate every image of
$M_{\mathbf b}\to M_{\mathbf a}$ with $\mathbf b\in L_q(\mathbf t)$ and
$\mathbf b\le\mathbf a$.  This gives the displayed formula for $\iota^*$.
The componentwise description of $\epsilon_V$ proves the assertion about
$\epsilon_V$.  Proposition~\ref{prop:recollement-ttf} gives the torsion pair
$(\operatorname{Im}\iota_*,\ker\iota^!)$, and the formula for $\iota^!$
identifies $\ker\iota^!=\operatorname{Im}j_*$.  Moreover,
$\operatorname{Im}\iota_*=\ker j^*$ is a Serre subcategory because $j^*$ is
exact.  Hence the torsion pair is hereditary.
\end{proof}


For $r\ge0$, put
\[
D_r(\mathbf t)=\{\mathbf a\in P_{\mathbf t}\mid\ell(\mathbf a)\le r\}.
\]

\begin{corollary}[Finite-box quotient formulas]
\label{cor:finite-quotient-formulas}
Let $\mathbf t\in\mathbb N^n$, let $1\le i\le n-1$, and put $q=n-i$.
Then
\[
\mathcal H_i(\mathbf t)/\mathcal H_{i-1}(\mathbf t)
\simeq\Rep(L_q(\mathbf t)).
\]
If $\mathbf t\ge\mathbf1$, then
\[
\Rep(L_q(\mathbf t))
\simeq
\prod_{\substack{F\subseteq[n]\\|F|=q}}
\Rep\left(\prod_{j\in F}[0,t_j-1]\right).
\]
Moreover, for $0\le j<n$ and $p=n-j$,
\begin{equation}\label{eq:finite-complementary-quotient}
\mathcal A^-_{\mathbf t}/\mathcal H_j(\mathbf t)
\simeq\Rep(D_{p-1}(\mathbf t)).
\end{equation}
\end{corollary}

\begin{proof}
The first equivalence is the Gabriel quotient supplied by
Proposition~\ref{prop:horizontal-embeddings}; equivalently, it is the
specialization of Theorem~\ref{thm:two-parameter}(3) with the box fixed.
When $\mathbf t\ge\mathbf1$, for $F\subseteq[n]$ with $|F|=q$, the mutually incomparable
components
\[
L_F(\mathbf t)=\{\mathbf a\in P_{\mathbf t}\mid
\operatorname{supp}(\mathbf a)=F\}
\]
are isomorphic to $\prod_{j\in F}[0,t_j-1]$ by subtracting $1$ in the
coordinates indexed by $F$.  This proves the product decomposition.  Finally,
\eqref{eq:finite-complementary-quotient} follows from
$P_{\mathbf t}=U_p(\mathbf t)\sqcup D_{p-1}(\mathbf t)$ and
Corollary~\ref{cor:order-ideal-consequences}, transported through the
finite-box equivalence of Proposition~\ref{prop:negative-box}.
\end{proof}

\begin{corollary}\label{cor:derived-horizontal-recollement}
Let $\mathbf t\in\mathbb N^n$, let $1\le i\le n-1$, and put $q=n-i$.
The abelian recollement in Proposition~\ref{prop:horizontal-embeddings} lifts
to the triangulated recollement:
\[
\begin{tikzcd}[column sep=3.4em]
{\scriptstyle D^b\!\left(\Rep(U_{q+1}(\mathbf t))\right)}
  \arrow[r, "{\scriptstyle \iota_*}" description]
& {\scriptstyle D^b\!\left(\Rep(U_q(\mathbf t))\right)}
  \arrow[l, bend right=20, "{\scriptstyle \mathbf L\iota^*}"']
  \arrow[l, bend left=20, "{\scriptstyle \iota^!}"]
  \arrow[r, "{\scriptstyle j^*}" description]
& {\scriptstyle D^b\!\left(\Rep(L_q(\mathbf t))\right)}
  \arrow[l, bend left=20, "{\scriptstyle j_!}"']
  \arrow[l, bend right=20, "{\scriptstyle j_*}"]
\end{tikzcd}.
\]
(the exact functors $\iota_*$, $\iota^!$, $j_!$, $j^*$, and $j_*$ are used degreewise).  Under the equivalences with the finite-box local cohomology categories, its three terms are
\[
D^b\!\left(\mathcal H_{i-1}(\mathbf t)\right),
\qquad
D^b\!\left(\mathcal H_i(\mathbf t)\right),
\qquad
D^b\!\left(\Rep(L_{n-i}(\mathbf t))\right),
\]
with the six functors transported from the displayed recollement.
Here $D^b(\mathcal H_{i-1}(\mathbf t))$ is identified with its fully faithful
image in $D^b(\mathcal H_i(\mathbf t))$ under the derived extension-by-zero
embedding.  Consequently
\[
\frac{D^b(\mathcal H_i(\mathbf t))}
     {D^b(\mathcal H_{i-1}(\mathbf t))}
\simeq
D^b\!\left(\Rep(L_{n-i}(\mathbf t))\right)
\simeq
D^b\!\left(
\frac{\mathcal H_i(\mathbf t)}{\mathcal H_{i-1}(\mathbf t)}
\right).
\]
If $\mathbf t\ge\mathbf1$, then moreover
\[
D^b\!\left(\Rep(L_q(\mathbf t))\right)
\simeq
\prod_{\substack{F\subseteq[n]\\|F|=q}}
D^b\!\left(\Rep\!\left(\prod_{j\in F}[0,t_j-1]\right)\right).
\]
In particular, for the squarefree box $\mathbf t=\mathbf1$,
\[
\frac{D^b(\mathcal H_i(\mathbf1))}
     {D^b(\mathcal H_{i-1}(\mathbf1))}
\simeq
\left(D^b(\Bbbk\text{-}\modcat)\right)^{\times\binom{n}{i}}.
\]
\end{corollary}

\begin{proof}
Since $U_q(\mathbf t)$ is finite,
Corollary~\ref{cor:fs-derived-recollement} gives the bounded-derived
recollement.  Lemma~\ref{lem:stratum-kan-exact} shows that in this particular
case $j_!$ and $j_*$ are already exact, so the displayed functors have the
stated form.  The displayed recollement induces the asserted Verdier quotient
equivalence, and the product and squarefree formulas are the derived forms of
Corollary~\ref{cor:finite-quotient-formulas}.
\end{proof}

Thus $j_*$ is extension by zero, whereas $j_!$ extends a rank-$q$ representation to all higher support ranks.  The next result gives a canonical and functorial description of this extension.

\subsection{The rank-layer resolution}\label{subsec:rank-resolution}

Continue to assume $1\le i\le n-1$ and $q=n-i$.
For $q\le r\le n$ define an exact functor
\[
E_r:\Rep(L_q(\mathbf t))\longrightarrow\Rep(L_r(\mathbf t))
\]
by
\[
(E_rV)_{\mathbf b}
=
\bigoplus_{\substack{F\subseteq\operatorname{supp}(\mathbf b)\\|F|=q}}
V_{\mathbf b_F}.
\]
If $\mathbf b\le\mathbf c$ lie in $L_r(\mathbf t)$, then they have the same support, and the structure map of $E_rV$ is the direct sum of the maps
$V_{\mathbf b_F}\to V_{\mathbf c_F}$.  Notice that
\[
E_q=\operatorname{id}_{\Rep(L_q(\mathbf t))}.
\]
Since $j_q=j$, put
\[
C_r(V)=j_{r,!}(E_rV);
\]
in particular, $C_q(V)=j_!V$.
At a vertex $\mathbf a$ with $A=\operatorname{supp}(\mathbf a)$ one has
\begin{equation}\label{eq:rank-term-stalk}
C_r(V)_{\mathbf a}
=
\bigoplus_{\substack{F\subseteq B\subseteq A\\|F|=q,\ |B|=r}}
V_{\mathbf a_F},
\end{equation}
with value zero when $r>|A|$.

For a finite set $X$, write $\binom{X}{m}$ for the set of $m$-element
subsets of $X$.
Fix the natural order on $[n]$.  For $v\in V_{\mathbf a_F}$, write
$v_{F,B}$ for its copy in the summand of \eqref{eq:rank-term-stalk} indexed by
$(F,B)$.  Define a natural transformation $d_r:C_r\to C_{r-1}$ by
\begin{equation}\label{eq:rank-differential}
d_r(v_{F,B})
=
\sum_{k\in B\setminus F}
(-1)^{\operatorname{pos}_{B\setminus F}(k)-1}
v_{F,B\setminus\{k\}},
\end{equation}
where $\operatorname{pos}_{B\setminus F}(k)$ is the position of $k$ in the increasing ordering of $B\setminus F$.  The identity on the coefficient $V_{\mathbf a_F}$ is understood in every summand.

\begin{theorem}[Functorial rank-layer resolution]\label{thm:rank-layer-resolution}
For every $V\in\Rep(L_q(\mathbf t))$, the maps \eqref{eq:rank-differential}, together with the canonical epimorphism
$\epsilon_V:j_!V\to j_*V$ of Proposition~\ref{prop:horizontal-embeddings}, form a functorial exact sequence
\begin{equation}\label{eq:rank-layer-resolution}
\boxed{
\begin{aligned}
0&\longrightarrow C_n(V)\xrightarrow{d_n}C_{n-1}(V)
\longrightarrow\cdots\longrightarrow C_{q+1}(V)\\
&\xrightarrow{d_{q+1}}C_q(V)=j_!V
\xrightarrow{\epsilon_V}j_*V\longrightarrow0.
\end{aligned}}
\end{equation}
\end{theorem}

\begin{proof}
The sign rule gives $d_{r-1}d_r=0$, since deleting two distinct elements in the two possible orders produces opposite signs.  Naturality follows because every component of $d_r$ is an identity on $V_{\mathbf a_F}$ and the structure maps of $V$ commute with these identities.

Exactness may be checked vertexwise.  Fix $\mathbf a\in U_q(\mathbf t)$, put $A=\operatorname{supp}(\mathbf a)$ and $s=|A|$, and first suppose $s>q$.  For a fixed $q$-subset $F\subseteq A$, the summands with coefficient $V_{\mathbf a_F}$ form
\[
0\longrightarrow
\Bbbk\!\left[\binom{A\setminus F}{s-q}\right]\longrightarrow\cdots\longrightarrow
\Bbbk\!\left[\binom{A\setminus F}{1}\right]\longrightarrow
\Bbbk\!\left[\{\varnothing\}\right]\longrightarrow0,
\]
tensored with $V_{\mathbf a_F}$.  With the differential \eqref{eq:rank-differential}, this is the augmented simplicial chain complex of the simplex on the nonempty vertex set $A\setminus F$, hence is exact.  Summing over all $F\in\binom{A}{q}$ proves exactness at $\mathbf a$.  If $s=q$, the only nonzero part is
\[
V_{\mathbf a}\xrightarrow{\mathrm{id}}V_{\mathbf a},
\]
coming from $C_q(V)_{\mathbf a}\to(j_*V)_{\mathbf a}$.  This proves the theorem.
\end{proof}

\begin{corollary}\label{cor:projective-rank-resolution}
If $V$ is projective in $\Rep(L_q(\mathbf t))$, then $E_rV$ is projective in
$\Rep(L_r(\mathbf t))$ and $C_r(V)$ is projective in
$\Rep(U_q(\mathbf t))$ for every $q\le r\le n$.  Consequently
\[
0\longrightarrow C_n(V)\longrightarrow\cdots\longrightarrow C_q(V)=j_!V
\xrightarrow{\epsilon_V}j_*V\longrightarrow0
\]
is a projective resolution of $j_*V$ of length at most $n-q=i$.  In the
Grothendieck group,
\[
[j_*V]=\sum_{r=q}^{n}(-1)^{r-q}[C_r(V)].
\]
\end{corollary}

\begin{proof}
Write
\[
L_q(\mathbf t)=\bigsqcup_{\substack{F\subseteq[n]\\|F|=q}}L_F(\mathbf t),
\qquad
L_r(\mathbf t)=\bigsqcup_{\substack{B\subseteq[n]\\|B|=r}}L_B(\mathbf t).
\]
For $F\subseteq[n]$ with $|F|=q$, write $V_F=V|_{L_F(\mathbf t)}$.  If $V$ is projective, then
each $V_F$ is a finite direct sum of principal projectives on
$L_F(\mathbf t)$.  For $F\subseteq B$, pullback along
the coordinate projection
\[
p_{B,F}:L_B(\mathbf t)\longrightarrow L_F(\mathbf t)
\]
sends the principal projective at $\mathbf a\in L_F(\mathbf t)$ to the
principal projective at the point of $L_B(\mathbf t)$ whose $F$-coordinates
are those of $\mathbf a$ and whose coordinates in $B\setminus F$ are all
$1$.  Moreover,
\[
(E_rV)|_{L_B(\mathbf t)}
\cong
\bigoplus_{\substack{F\subseteq B\\|F|=q}}p_{B,F}^*V_F.
\]
Hence each pullback $p_{B,F}^*V_F$ is projective, and therefore so is
$E_rV$.  Finally, $j_{r,!}$ is left adjoint to the exact restriction
functor $j_r^*$, so it preserves projectives.  Thus
$C_r(V)=(j_{r,!}\circ E_r)(V)$ is projective.  The remaining assertions follow from
Theorem~\ref{thm:rank-layer-resolution}.
\end{proof}

\begin{remark}\label{rem:dual-rank-resolution}
Order reversal gives a dual co-rank coresolution.  Put
\[
\ell^\vee(\mathbf a)=|\{k\in[n]\mid a_k<t_k\}|,
\qquad
U_q^\vee(\mathbf t)=\{\mathbf a\in P_{\mathbf t}\mid
\ell^\vee(\mathbf a)\ge q\},
\]
and let $L_r^\vee(\mathbf t)$ be the stratum on which
$\ell^\vee(\mathbf a)=r$.  The order-reversing bijection
$\tau_{\mathbf t}(\mathbf a)=\mathbf t-\mathbf a$ identifies
$U_q^\vee(\mathbf t)$ with $U_q(\mathbf t)^{\mathrm{op}}$ and
$L_r^\vee(\mathbf t)$ with $L_r(\mathbf t)^{\mathrm{op}}$.  Consequently the
exact contravariant equivalence $\tau_{\mathbf t}^*\circ D$ sends the
rank-layer resolution of Theorem~\ref{thm:rank-layer-resolution} in
$\Rep(U_q(\mathbf t))$ to a right-Kan coresolution in
$\Rep(U_q^\vee(\mathbf t))$, whose terms are induced from the co-rank strata
$L_r^\vee(\mathbf t)$.  We shall not need its explicit formula below.
\end{remark}

\subsection{Nakayama--Serre exchange for rank-layer resolutions}\label{subsec:higher-nakayama}

Retain the notation and assumptions of the preceding subsection.  For a
finite poset $P$, write
\[
\boldsymbol\nu_P
=
DA_P\otimes_{A_P}^{\mathbf L}-
:
D^b(\Rep(P))\longrightarrow D^b(\Rep(P))
\]
for the derived Nakayama functor.  With the incidence-algebra convention fixed
in Section~\ref{sec:homological}, one has $A_{P_{\mathbf t}}=\Lambda_{\mathbf t}$.
Since $A_P$ has finite global dimension, $\boldsymbol\nu_P$ is the Serre
functor of $D^b(\Rep(P))$ and hence agrees, up to canonical natural
isomorphism, with the right Serre functor $\mathbb S_P^R$ introduced above.
Thus $\boldsymbol\nu_P\simeq\mathbb S_P^R$.  For $X,Y\in D^b(\Rep(P))$
there are bifunctorial isomorphisms
\[
D\Hom_{D^b(\Rep(P))}(X,Y)
\cong
\Hom_{D^b(\Rep(P))}(Y,\boldsymbol\nu_P X).
\]
It is therefore an autoequivalence, with quasi-inverse
\[
\boldsymbol\nu_P^{-1}
\simeq
\mathbf R\Hom_{A_P}(DA_P,-).
\]

For $q\le r\le n$, recall the finite-box inclusion
\[
j_r:L_r(\mathbf t)\hookrightarrow U_q(\mathbf t).
\]
By Lemma~\ref{lem:stratum-kan-exact}, both $j_{r,!}$ and $j_{r,*}$ are exact;
we use the same symbols for the induced triangle functors on bounded derived
categories.

\begin{proposition}[Finite-box Nakayama--Serre exchange]\label{prop:nakayama-kan-exchange}
For every $q\le r\le n$, there are canonical natural isomorphisms
\begin{equation}\label{eq:nakayama-kan-exchange}
\boldsymbol\nu_{U_q(\mathbf t)}\circ j_{r,!}
\xRightarrow{\ \sim\ }
j_{r,*}\circ\boldsymbol\nu_{L_r(\mathbf t)},
\qquad
\boldsymbol\nu_{U_q(\mathbf t)}^{-1}\circ j_{r,*}
\xRightarrow{\ \sim\ }
j_{r,!}\circ\boldsymbol\nu_{L_r(\mathbf t)}^{-1}.
\end{equation}
Both sides are triangle functors
\[
D^b\!\left(\Rep(L_r(\mathbf t))\right)
\longrightarrow
D^b\!\left(\Rep(U_q(\mathbf t))\right).
\]
Thus Nakayama--Serre duality exchanges the left-Kan and right-Kan extensions
from each support-rank stratum.
\end{proposition}

\begin{proof}
Apply Theorem~\ref{thm:right-serre-exchange} to
$j_r:L_r(\mathbf t)\hookrightarrow U_q(\mathbf t)$.  Both posets are finite,
so Proposition~\ref{prop:right-serre-existence} identifies their right Serre
functors with the derived Nakayama functors.  Since $j_{r,!}$ and $j_{r,*}$
are exact, their derived functors are represented by the same Kan extension
functors.  This gives the first isomorphism; the second follows by composing
with the inverse Serre autoequivalences.
\end{proof}

If $V$ is projective in $\Rep(L_q(\mathbf t))$, then Corollary~\ref{cor:projective-rank-resolution}
shows that $E_rV$ and $C_r(V)=j_{r,!}(E_rV)$ are projective for every
$q\le r\le n$.  Hence Proposition~\ref{prop:nakayama-kan-exchange} applies
term by term to the rank-layer resolution:
\[
\boldsymbol\nu_{U_q(\mathbf t)}\bigl(C_r(V)\bigr)
\cong
j_{r,*}\!\left(\boldsymbol\nu_{L_r(\mathbf t)}(E_rV)\right).
\]
Thus the left-Kan rank resolution is transformed functorially into a
right-Kan, or costandard, rank complex.

\medskip\noindent\textbf{Relation with Brun--Fl\o ystad.}
We use the same symbol $\pi_{\mathbf t}$ for the equivalence induced on
bounded derived categories.  On the full box $P_{\mathbf t}$,
the natural quasi-isomorphism used in the proof of
Proposition~\ref{prop:naka-ext}, together with
\cite[Proposition~3.9]{BF}, gives
\begin{equation}\label{eq:BF-derived-nakayama-comparison}
\boldsymbol\nu_{P_{\mathbf t}}\circ\pi_{\mathbf t}
\simeq
\bigl(\pi_{\mathbf t}\circ\mathcal N_{\mathbf t}^{\mathbf1}\bigr)[n].
\end{equation}
Here we use the standard natural identification, valid for every perfect
complex $X$ of left $\Lambda_{\mathbf t}$-modules,
\[
D\mathbf R\!\Hom_{\Lambda_{\mathbf t}}(X,\Lambda_{\mathbf t})
\simeq
D\Lambda_{\mathbf t}\otimes_{\Lambda_{\mathbf t}}^{\mathbf L}X.
\]
Thus the full-box derived Nakayama functor is, up to a uniform shift, the
incidence-algebra realization of the first full Nakayama construction of
Brun--Fl\o ystad.  After fixing $i<n$ and $q=n-i$, however,
$\boldsymbol\nu_{U_q(\mathbf t)}$ is the intrinsic Serre functor of the
smaller incidence algebra $A_{U_q(\mathbf t)}$ and governs the internal
homological structure of the already formed category
$\mathcal H_i(\mathbf t)$.  It should not be interpreted as an additional BF
iterate or as another local cohomology functor.

\medskip\noindent\textbf{The costandard rank complex.}
The preceding proposition now gives a direct derived description of the
Nakayama image of the canonical right section $j_*V$.


\begin{theorem}[Rank complex for cohomological Nakayama images]\label{thm:higher-nakayama-rank-complex}
Let $V$ be projective in $\Rep(L_q(\mathbf t))$.  For $q\le r\le n$, the
object $E_rV$ is projective in $\Rep(L_r(\mathbf t))$; hence
$\boldsymbol\nu_{L_r(\mathbf t)}(E_rV)$ is concentrated in degree zero, and
\[
j_{r,*}\!\left(\boldsymbol\nu_{L_r(\mathbf t)}(E_rV)\right)
\in\Rep(U_q(\mathbf t)).
\]
Place this object in cohomological degree $q-r$.  Then there is a functorial
isomorphism in $D^b(\Rep(U_q(\mathbf t)))$
\begin{equation}\label{eq:higher-nakayama-complex}
\boxed{
\boldsymbol\nu_{U_q(\mathbf t)}(j_*V)
\simeq
\left[
\begin{aligned}
\cdots&\longrightarrow 0\longrightarrow
j_{n,*}\!\left(\boldsymbol\nu_{L_n(\mathbf t)}(E_nV)\right)\\
&\longrightarrow\cdots\longrightarrow
j_{q+1,*}\!\left(\boldsymbol\nu_{L_{q+1}(\mathbf t)}(E_{q+1}V)\right)\\
&\longrightarrow
j_*\!\left(\boldsymbol\nu_{L_q(\mathbf t)}(V)\right)
\longrightarrow 0\longrightarrow\cdots
\end{aligned}
\right].}
\end{equation}
The displayed complex is zero outside cohomological degrees
$q-n,\ldots,0$.  Consequently, for $0\le p\le n-q=i$,
\begin{equation}\label{eq:higher-nakayama-ext}
H^{-p}\!\left(\boldsymbol\nu_{U_q(\mathbf t)}(j_*V)\right)
\cong
D\Ext^p_{A_{U_q(\mathbf t)}}
\!\left(j_*V,A_{U_q(\mathbf t)}\right),
\end{equation}
as objects of $\Rep(U_q(\mathbf t))$, and these cohomology objects vanish for
$p>i$.  The complex
\eqref{eq:higher-nakayama-complex} is the costandard rank complex associated
with $V$.
\end{theorem}

\begin{proof}
By Corollary~\ref{cor:projective-rank-resolution},
\eqref{eq:rank-layer-resolution} is a projective resolution of $j_*V$ of
length at most $i$.  Since every $C_r(V)$ is projective, the derived Nakayama
functor is computed by applying it term by term.  Moreover $E_rV$ is
projective, and Proposition~\ref{prop:nakayama-kan-exchange} gives
\[
\boldsymbol\nu_{U_q(\mathbf t)}\bigl(C_r(V)\bigr)
=
\boldsymbol\nu_{U_q(\mathbf t)}\bigl(j_{r,!}(E_rV)\bigr)
\simeq
j_{r,*}\!\left(\boldsymbol\nu_{L_r(\mathbf t)}(E_rV)\right).
\]
Naturality of the exchange identifies these terms compatibly with the
rank-layer differentials, yielding \eqref{eq:higher-nakayama-complex}.
Finally, for a finite-dimensional algebra $A$ and a bounded projective
resolution of $M$,
\[
H^{-p}(DA\otimes_A^{\mathbf L}M)
\cong
D\Ext_A^p(M,A).
\]
Applying this with $A=A_{U_q(\mathbf t)}$ gives
\eqref{eq:higher-nakayama-ext}; the vanishing for $p>i$ follows from the
length of the projective resolution.
\end{proof}

\begin{thebibliography}{99}
\bibitem[ARS]{ARS} M.~Auslander, I.~Reiten and S.~O.~Smal\o, \newblock \emph{Representation Theory of Artin Algebras}, \newblock Cambridge Studies in Advanced Mathematics, vol.~36, Cambridge University Press, Cambridge, 1995.
\bibitem[ARsyzygy]{ARsyzygy} M.~Auslander and I.~Reiten, \newblock $k$-Gorenstein algebras and syzygy modules, \newblock \emph{J. Pure Appl. Algebra} \textbf{92} (1994), no.~1, 1--27. \newblock \href{https://doi.org/10.1016/0022-4049(94)90044-2}{doi:10.1016/0022-4049(94)90044-2}.
\bibitem[BBD]{BBD} A.~A.~Beilinson, J.~Bernstein and P.~Deligne, \newblock Faisceaux pervers, \newblock \emph{Ast\'erisque} \textbf{100} (1982), 5--171.
\bibitem[BF]{BF} M.~Brun and G.~Fl\o ystad, \newblock The Auslander--Reiten translate on monomial rings, \newblock \emph{Adv. Math.} \textbf{226} (2011), no.~1, 952--991. \newblock \href{https://doi.org/10.1016/j.aim.2010.07.006}{doi:10.1016/j.aim.2010.07.006}.
\bibitem[BH]{BH} W.~Bruns and J.~Herzog, \newblock \emph{Cohen--Macaulay Rings}, \newblock revised ed., Cambridge Studies in Advanced Mathematics, vol.~39, Cambridge University Press, Cambridge, 1998.
\bibitem[FGR]{FGR} R.~M.~Fossum, P.~A.~Griffith and I.~Reiten, \newblock \emph{Trivial Extensions of Abelian Categories}, \newblock Lecture Notes in Mathematics, vol.~456, Springer-Verlag, Berlin--Heidelberg--New York, 1975.
\bibitem[FP]{FP} V.~Franjou and T.~Pirashvili, \newblock Comparison of abelian categories recollements, \newblock \emph{Doc. Math.} \textbf{9} (2004), 41--56. \newblock \href{https://doi.org/10.4171/DM/156}{doi:10.4171/DM/156}.
\bibitem[Gabriel]{Gabriel} P.~Gabriel, \newblock Des cat\'egories ab\'eliennes, \newblock \emph{Bull. Soc. Math. France} \textbf{90} (1962), 323--448.
\bibitem[IM]{IM} O.~Iyama and R.~Marczinzik, \newblock Distributive lattices and Auslander regular algebras, \newblock \emph{Adv. Math.} \textbf{398} (2022), 108233. \newblock \href{https://doi.org/10.1016/j.aim.2022.108233}{doi:10.1016/j.aim.2022.108233}.
\bibitem[Miller]{Miller} E.~Miller, \newblock The Alexander duality functors and local duality with monomial support, \newblock \emph{J. Algebra} \textbf{231} (2000), no.~1, 180--234. \newblock \href{https://doi.org/10.1006/jabr.2000.8359}{doi:10.1006/jabr.2000.8359}.
\bibitem[PV]{PV} C.~Psaroudakis and J.~Vit\'oria, \newblock Recollements of module categories, \newblock \emph{Appl. Categ. Structures} \textbf{22} (2014), no.~4, 579--593. \newblock \href{https://doi.org/10.1007/s10485-013-9323-x}{doi:10.1007/s10485-013-9323-x}.
\bibitem[Psaroudakis]{Psaroudakis} C.~Psaroudakis, \newblock Homological theory of recollements of abelian categories, \newblock \emph{J. Algebra} \textbf{398} (2014), 63--110. \newblock \href{https://doi.org/10.1016/j.jalgebra.2013.09.020}{doi:10.1016/j.jalgebra.2013.09.020}.
\end{thebibliography}
\end{document}
